% =====================================================================
%  U10 —— 第 7 章：交织算子
%  底本：E. Frenkel, Langlands Correspondence for Loop Groups (2007)
%  PDF p209–p231（书页编号见页眉：197–219）
% =====================================================================

\chapter{交织算子}

我们现在来证明定理 4.3.2。证明分本章与下一章给出，依据 [Frenkel 2005b]。该定理最初由 [Feigin and Frenkel 1992; Frenkel 1991] 证明（我们指出，一个密切相关的命题，即定理 8.3.1，是由 V. Drinfeld 猜想的）。

本章开头先概述证明的整体策略（见 §7.1），然后完成证明的前两步。接着我们发展 Wakimoto 模之间交织算子的理论，这是定理 4.3.2 证明中的重要一步。先在 §7.2 对 $\widehat{\mathfrak{sl}}_2$ 的情形来做：我们显式构造交织算子，称之为第一类与第二类屏蔽算子。在 §7.3 中我们利用这些算子与抛物诱导函子，对任意仿射 Kac--Moody 代数构造 Wakimoto 模之间的交织算子。

% [Tutor Note: 本章与第 8 章合起来构成定理 4.3.2（中心 ≅ opers 上的函数）的证明，依据 [Frenkel 2005b]。读者宜先记住第 197 页（PDF p.209）列出的六步路线图：本章只做第 1–4 步（单射性、中心落入 π₀、构造屏蔽算子、中心落入核之交），第 5–6 步（特征标比较与 Miura opers）留到第 8 章。这是全书唯一一处“先给路线图再分两章实施”的写法。]

\begin{keypoints}
    \item \textbf{核心思想：}把中心的研究化归为在 Wakimoto 模的自由场实现 $W_{0,\kappa_c}=M_{\mathfrak{g}}\otimes\pi_0$ 中，把 $V_{\kappa_c}(\mathfrak{g})$ 与 $\mathfrak{z}(\hat{\mathfrak{g}})$ 刻画成屏蔽算子核之交。
    \item \textbf{数学工具：}模上的顶点算子与 OPE、Friedan--Martinec--Shenker 玻色化、第一类/第二类屏蔽算子、半无限抛物诱导。
    \item \textbf{关键结论：}$w_\kappa:V_\kappa(\mathfrak{g})\hookrightarrow W_{0,\kappa}$ 是单射；$\mathfrak{z}(\hat{\mathfrak{g}})$ 落在 $\bigcap_i\operatorname{Ker}\overline{V}_i[1]$ 中（命题 7.3.6）。
    \item \textbf{记号提醒：}$w_\kappa$ 是自由场实现同态；$S_k,\widetilde{S}_k$ 是 $\mathfrak{sl}_2$ 的第一类、第二类屏蔽算子；带上横线的 $\overline{S}_i,\overline{V}_i[1]$ 是其在临界水平的极限。
\end{keypoints}

\section{证明的策略}

我们证明定理 4.3.2 的策略如下。在定理 6.1.6 中我们构造了顶点代数的自由场实现同态
\transnote{卡点：本章目标定理在原文中时而称“定理 4.3.2”、时而称“定理 6.1.6”，同一段内两种编号并存，读者易以为漏了一个定理。背景：定理 4.3.2 断言中心 $\mathfrak{z}(\hat{\mathfrak{g}})$ 与 opers 空间上的函数代数同构；定理 6.1.6 是同一结论经自由场实现（Wakimoto 模）表述的版本，本章与第 8 章合起来证明的是后者，再由等价性回到前者。指代落点：§4.3 定理 4.3.2、§6.1 定理 6.1.6；原书 p.197 章首用 4.3.2，p.198--199 三处用 6.1.6。译文照原文保留两种编号，未作统一。[依据：原书 p.197--199 各处上下文][来源：本书 p.197。]}
\begin{equation}
    w_{\kappa_c}:V_{\kappa_c}(\mathfrak{g})\longrightarrow M_{\mathfrak{g}}\otimes\pi_0.
    \tag{7.1-1}
\end{equation}

\textbf{第 1 步。}我们将证明同态 (7.1-1) 是单射。

\textbf{第 2 步。}我们将证明 $\mathfrak{z}(\hat{\mathfrak{g}})\subset V_{\kappa_c}(\mathfrak{g})$ 在 $w_{\kappa_c}$ 下的像包含在 $\pi_0\subset M_{\mathfrak{g}}\otimes\pi_0$ 中。

于是，我们需要描述 $\mathfrak{z}(\hat{\mathfrak{g}})$ 在 $\pi_0$ 中的像。

\textbf{第 3 步。}我们将构造屏蔽算子 $\overline{S}_i,\,i=1,\dots,\ell$，从 $W_{0,\kappa_c}=M_{\mathfrak{g}}\otimes\pi_0$ 到某些其他模，它们与 $\widehat{\mathfrak{g}}_{\kappa_c}$ 的作用交换。

\textbf{第 4 步。}我们将证明 $V_{\kappa_c}(\mathfrak{g})$ 在 $w_{\kappa_c}$ 下的像包含在
\[
    \bigcap_{i=1}^{\ell}\operatorname{Ker}\overline{S}_i
\]
中。由此推出 $\mathfrak{z}(\hat{\mathfrak{g}})$ 的像包含在
\[
    \bigcap_{i=1}^{\ell}\operatorname{Ker}\overline{V}_i[1]
\]
中，其中 $\overline{V}_i[1]$ 是 $\overline{S}_i$ 到 $\pi_0$ 的限制。

\textbf{第 5 步。}利用我们这些模的分次伴随以及命题 6.3.3 中的同构 $W^+_{0,\kappa_c}\simeq\mathbb{M}_{0,\kappa_c}$，我们将求出 $\mathfrak{z}(\hat{\mathfrak{g}})$ 的特征标。我们将证明它等于 $\bigcap_{i=1}^{\ell}\operatorname{Ker}\overline{V}_i[1]$ 的特征标。于是得到
\[
    \mathfrak{z}(\hat{\mathfrak{g}})=\bigcap_{i=1}^{\ell}\operatorname{Ker}\overline{V}_i[1].
\]

\textbf{第 6 步。}利用 Miura opers，我们将证明存在自然同构
\[
    \operatorname{Fun}\Op_{{}^LG}(D)\simeq\bigcap_{i=1}^{\ell}\operatorname{Ker}\overline{V}_i[1].
\]
于是得到
\[
    \mathfrak{z}(\hat{\mathfrak{g}})\simeq\operatorname{Fun}\Op_{{}^LG}(D).
\]

这将给出定理 6.1.6 的证明，因为我们将证明上述等同保持群 $\operatorname{Aut}\mathcal{O}$ 的自然作用。我们还将证明这个同构满足各种其他相容性。

本书给出的定理 6.1.6 的证明紧跟 [Frenkel 2005b]。此证明与 [Feigin and Frenkel 1992; Frenkel 1991] 中的原始证明有两点不同。首先，我们使用第二类而非第一类屏蔽算子，它们的 $\kappa\to\kappa_c$ 极限更容易研究。其次，我们利用临界水平最高权为 $0$ 的 Verma 模与某个 Wakimoto 模之间的同构，并计算奇异向量空间的分次伴随，来估计中心的特征标。

下面我们执行上述计划中的第 1、2 步。第 3、4 步将在本章余下部分处理。

% [Tutor Note: 原文在此把目标定理时而称为定理 4.3.2、时而称为定理 6.1.6（章首一段说要证明定理 4.3.2，而第 197–198 页三处说定理 6.1.6）。两处指同一结论的不同表述：定理 4.3.2 是中心与 Fun Op 的同构，定理 6.1.6 是其经由自由场实现的等价形式。译文照原文保留这两种编号。]

\begin{keypoints}
    \item \textbf{核心思想：}证明分六步；本章完成第 1、2 步，并把第 3、4 步（屏蔽算子）在本章后半完成。
    \item \textbf{数学工具：}PBW 滤过与分次伴随、射流概形、优势态射、Grothendieck alteration、张量积分解。
    \item \textbf{关键结论：}$w_\kappa$ 对任意 $\kappa$ 是单射（命题 7.1.1）；$\mathfrak{z}(\hat{\mathfrak{g}})$ 的像落在 $\pi_0$ 中（引理 7.1.2）。
    \item \textbf{记号提醒：}$\mathfrak{b}_-$ 是下 Borel、$\mathfrak{b}_+$ 是上 Borel；关键态射是 $p:\mathfrak{b}_-\times N_+\to\mathfrak{g}$。
\end{keypoints}

\subsection{有限维情形}

先从第 1 步的陈述开始。事实上，我们将证明一个更一般的结果，它对任意 $\kappa$ 成立，而不只对 $\kappa_c$。

在给出该陈述的证明之前，先看它在有限维情形的类比是有启发性的。

回忆 §5.2 中我们构造了同态 (5.2-10)，
\begin{equation}
    \widetilde{\rho}:U(\mathfrak{g})\longrightarrow\operatorname{Fun}\mathfrak{h}^*\otimes_{\mathbb{C}}\mathcal{D}(N_+).
    \tag{7.1-2}
\end{equation}
我们希望证明这个同态是单射。

考虑 $U(\mathfrak{g})$ 上的 Poincaré--Birkhoff--Witt 滤过，以及 $\operatorname{Fun}\mathfrak{h}^*\otimes_{\mathbb{C}}\mathcal{D}(N_+)$ 上如下定义的滤过：其第 $n$ 项是直和
\[
    \bigoplus_{i=0}^{n}(\operatorname{Fun}\mathfrak{h}^*)_{\le i}\otimes\mathcal{D}(N_+)_{\le(n-i)}.
\]
这里 $(\operatorname{Fun}\mathfrak{h}^*)_{\le i}$ 表示次数不超过 $i$ 的多项式空间，$\mathcal{D}(N_+)_{\le(n-i)}$ 是阶数不超过 $n-i$ 的微分算子空间。

由 §5.2 中给出的公式显然可知，同态 $\widetilde{\rho}$ 保持这些滤过。考虑相应的分次伴随代数同态。回忆 $\operatorname{gr}U(\mathfrak{g})=\operatorname{Fun}\mathfrak{g}^*$ 与 $\operatorname{gr}\mathcal{D}(N_+)=\operatorname{Fun}T^*N_+$。显然 $\operatorname{gr}\operatorname{Fun}\mathfrak{h}^*=\operatorname{Fun}\mathfrak{h}^*$。于是我们得到交换代数的同态
\begin{equation}
    \operatorname{Fun}\mathfrak{g}^*\longrightarrow\operatorname{Fun}\mathfrak{h}^*\otimes_{\mathbb{C}}\operatorname{Fun}T^*N_+.
    \tag{7.1-3}
\end{equation}
它对应仿射代数簇的态射
\[
    \mathfrak{h}^*\times T^*N_+\longrightarrow\mathfrak{g}^*.
\]
注意 $N_+$ 的余切丛等同于 $N_+$ 上以 $\mathfrak{n}_+$ 为纤维的平凡向量丛。由于 $\mathfrak{b}_+=\mathfrak{h}\oplus\mathfrak{n}$，可得 $\mathfrak{h}^*\oplus\mathfrak{n}_+^*=\mathfrak{b}_+^*$。\transnote{卡点：原文此处印作 $\mathfrak{b}_+=\mathfrak{h}\oplus\mathfrak{n}$，同句结论却用 $\mathfrak{n}_+$（$\mathfrak{h}^*\oplus\mathfrak{n}_+^*=\mathfrak{b}_+^*$）；照 $\mathfrak{n}$ 读，等式不成立。背景：上 Borel 的分解是 $\mathfrak{b}_+=\mathfrak{h}\oplus\mathfrak{n}_+$，其中 $\mathfrak{n}_+=\bigoplus_{\alpha>0}\mathfrak{g}_\alpha$；$N_+$ 的余切丛纤维为 $\mathfrak{n}_+$，对偶后正给出 $\mathfrak{b}_+^*=\mathfrak{h}^*\oplus\mathfrak{n}_+^*$。指代落点：§7.1.1，(7.1-3) 的态射 $\mathfrak{h}^*\times T^*N_+\to\mathfrak{g}^*$，原书 p.200。此处 $\mathfrak{n}$ 当是 $\mathfrak{n}_+$ 的排印脱漏，译文照原文保留。[依据：原书 p.200 同句两处记号自相参照][来源：本书 p.199。]}
因此上述态射等价于态射
\[
    \mathfrak{b}_+^*\times N_+\longrightarrow\mathfrak{g}^*.
\]
利用 $\mathfrak{g}$ 上非退化不变内积 $\kappa_0$，我们得到态射
\begin{equation}
    p:\mathfrak{b}_-\times N_+\longrightarrow\mathfrak{g}.
    \tag{7.1-4}
\end{equation}
后一个态射容易显式描述。由定义可知它把
\[
    (x,g)\in\mathfrak{b}_-\times N_+\longmapsto gxg^{-1}\in\mathfrak{g}.
\]
因此 $p$ 在 $\mathfrak{g}$ 中的像由所有具备下述性质的元素组成：它们属于某个 Borel 子代数 $\mathfrak{b}$，而该 $\mathfrak{b}$ 对应的 $\Fl=G/B_-$ 中的点位于开的 $N_+$-轨道 $\mathcal{U}=N_+\cdot[1]$ 中（我们说这样的 $\mathfrak{b}$ 与 $\mathfrak{b}_-$ 处于\textbf{一般相对位置}）。因此 $p$ 的像是 $\mathfrak{g}$ 中的开稠密子集。换言之，$p$ 是优势态射。此外，像中的一般元素只属于唯一一个这样的 Borel 子代数 $\mathfrak{b}$，故 $p$ 一般地一一。于是同态 (7.1-3) 是单射，从而 (7.1-2) 中的同态 $\widetilde{\rho}$ 也是单射。
\transnote{卡点：为什么“$p$ 优势且一般地一一”能给出 (7.1-3) 单射、进而 $w_\kappa$ 单射，原文一笔带过。背景：对代数闭域上的仿射簇，态射优势（像稠密）当且仅当其拉回映射在函数环上单射——这正是 (7.1-3) 单射的来源；“一般地一一”则用于保证分次伴随层面的单射能沿滤过传回 $w_\kappa$ 本身。严格论证要在 $\mathfrak{g}^*$ 的一般点处取局部环，并先经 $\kappa_0$ 把 $\mathfrak{g}$ 与 $\mathfrak{g}^*$ 等同。所谓“一般相对位置”，指 $\mathfrak{b}$ 落在 $\Fl$ 的开 $N_+$-轨道 $\mathcal{U}$ 中，等价于 $\mathfrak{b}$ 与 $\mathfrak{b}_-$ 横截。指代落点：§7.1.1，(7.1-3)--(7.1-4)，原书 p.200--201。[来源：经典教材 Hartshorne, \textit{Algebraic Geometry}, GTM 52, 第 II 章（优势态射与函数环）；原书 p.200--201]}

注意态射 $p$ 可以在 \textbf{Grothendieck alteration} 的框架下重新表述。设 $\widetilde{\mathfrak{g}}$ 是由点对 $(\mathfrak{b},x)$ 构成的簇，其中 $\mathfrak{b}$ 是 $\mathfrak{g}$ 中的 Borel 子代数，$x\in\mathfrak{b}$。有自然态射 $\widetilde{\mathfrak{g}}\to\Fl$，把 $(\mathfrak{b},x)$ 映到 $\mathfrak{b}\in\Fl$。此映射把 $\widetilde{\mathfrak{g}}$ 等同于旗簇 $\Fl$ 上的向量丛，其在 $\mathfrak{b}\in\Fl$ 上的纤维是向量空间 $\mathfrak{b}$。还有态射 $\widetilde{\mathfrak{g}}\to\mathfrak{g}$，把 $(\mathfrak{b},x)$ 映到 $x$。

现在设 $\mathcal{U}$ 是大胞腔，即 $\Fl$ 的开 $N_+$-轨道，$\widetilde{\mathcal{U}}$ 是它在 $\widetilde{\mathfrak{g}}$ 中的原像。换言之，$\widetilde{\mathcal{U}}$ 由那些点对 $(\mathfrak{b},x)$ 组成，其中 $\mathfrak{b}$ 与 $\mathfrak{b}_-$ 处于一般相对位置。则 $\widetilde{\mathcal{U}}$ 自然同构于 $\mathfrak{b}_-\times N_+$，而态射 $\widetilde{\mathfrak{g}}\to\mathfrak{g}$ 在 $\widetilde{\mathcal{U}}$ 上的限制就是上述态射 $p$。

% [Tutor Note: “$p$ 优势且一般地一一，故函数环同态是单射”这一步，本质上是代数几何中“优势态射诱导的函数环同态为单射”的断言（对整簇而言），其严格证明要在 generic point 处取局部环。原文在此一笔带过；另外 $(x,g)\mapsto gxg^{-1}$ 把 $\mathfrak{b}_-\times N_+$ 送到 $\mathfrak{g}$，其像恰是那些落在与 $\mathfrak{b}_-$ 横截的 Borel 中的元素之并，这正是 Grothendieck alteration 的“大胞腔”部分。]

\subsection{单射性}

现在可以来考虑仿射李代数的类似问题，它对应我们计划的第 1 步。

\begin{proposition}
    定理 6.2.1 中的同态 $w_\kappa$ 对任意 $\kappa$ 都是单射。
\end{proposition}

\begin{proof}
我们沿用有限维情形中的论证。即在 $V_\kappa(\mathfrak{g})$ 与 $W_{0,\kappa}=M_{\mathfrak{g}}\otimes\pi_0^{\kappa-\kappa_c}$ 上引入被 $w_\kappa$ 保持的滤过，然后证明诱导映射 $\operatorname{gr}w_\kappa:\operatorname{gr}V_\kappa(\mathfrak{g})\to\operatorname{gr}W_{0,\kappa}$（它其实与 $\kappa$ 无关）是单射。

$U(\hat{\mathfrak{g}})$ 上的 Poincaré--Birkhoff--Witt 滤过按 §2.2.5 所述的方式在 $V_\kappa(\mathfrak{g})$ 上诱导滤过。

现在在 $W_{0,\kappa}$ 上定义滤过 $\{W^{\le p}_{0,\kappa}\}$。按定义，$W^{\le p}_{0,\kappa}$ 是由 $a_{\alpha,n}$、$a^*_{\alpha,n}$ 与 $b_{i,n}$ 的单项式张成的空间，这些单项式在 $a_{\alpha,n}$ 与 $b_{i,n}$ 中的合次数不超过 $p$（这类似于上面有限维情形所用的滤过）。由同态 $w_\kappa$ 的构造显然可知它保持这些滤过。此外，这些是顶点代数的滤过，其分次伴随空间是交换顶点代数，从而是交换代数。因此 $w_\kappa$ 的分次伴随 $\operatorname{gr}w_\kappa$ 是这些交换代数的同态。

现在描述相应的交换代数 $\operatorname{gr}V_\kappa(\mathfrak{g})$ 与 $\operatorname{gr}W_{0,\kappa}$，以及同态 $\operatorname{gr}w_\kappa:\operatorname{gr}V_\kappa(\mathfrak{g})\to\operatorname{gr}W_{0,\kappa}$。我们用非退化不变内积 $\kappa_0$ 把 $\mathfrak{g}$ 与 $\mathfrak{g}^*$ 等同。

设 $J(\mathfrak{b}_-\times N_+)=J\mathfrak{b}_-\times JN_+$ 与 $J\mathfrak{g}$ 分别是 $\mathfrak{b}_-\times N_+$ 与 $\mathfrak{g}$ 的无限射流概形，按 §3.4.2 定义。则 $\operatorname{gr}V_\kappa(\mathfrak{g})=\operatorname{Fun}J\mathfrak{g}$，$\operatorname{gr}W_{0,\kappa}=\operatorname{Fun}J(\mathfrak{b}_-\times N_+)$。同态
\begin{equation}
    \operatorname{gr}w_\kappa:\operatorname{Fun}J\mathfrak{g}\longrightarrow\operatorname{Fun}J(\mathfrak{b}_-\times N_+)
    \tag{7.1-5}
\end{equation}
对应射流概形的态射
\[
    J(\mathfrak{b}_-\times N_+)\longrightarrow J\mathfrak{g}.
\]
由构造显然可知，它就是由 $J$ 的函子性对应于公式 (7.1-4) 中态射 $p$ 的那个态射 $Jp$。

但 $p$ 是优势态射且一般地一一。设 $\mathfrak{g}_{\mathrm{gen}}$ 是 $p$ 在 $\mathfrak{g}$ 中的像里使 $p$ 为一一的轨迹。则 $\mathfrak{g}_{\mathrm{gen}}$ 在 $\mathfrak{g}$ 中开稠密。因此 $J\mathfrak{g}_{\mathrm{gen}}$ 在 $J\mathfrak{g}$ 中开稠密，而 $J\mathfrak{g}_{\mathrm{gen}}$ 显然在 $Jp$ 的像中。于是 $Jp$ 是优势态射且一般地一一。因此函数环的同态 (7.1-5) 是单射。这也推出 $w_\kappa$ 是单射。
\end{proof}

现在专门考虑临界水平。顶点代数 $V_{\kappa_c}(\mathfrak{g})$ 含有交换子代数 $\mathfrak{z}(\hat{\mathfrak{g}})$，即其中心。回忆 $\mathfrak{z}(\hat{\mathfrak{g}})$ 是 $V_{\kappa_c}(\mathfrak{g})$ 中 $\mathfrak{g}[[t]]$-不变向量组成的空间。另一方面，$W_{0,\kappa_c}=M_{\mathfrak{g}}\otimes\pi_0$ 含有交换子代数 $\pi_0$，它是其中心。

\begin{lemma}
    $\mathfrak{z}(\hat{\mathfrak{g}})\subset V_{\kappa_c}(\mathfrak{g})$ 在 $w_{\kappa_c}$ 下于 $W_{0,\kappa_c}$ 中的像包含在 $\pi_0\subset W_{0,\kappa_c}$ 中。
\end{lemma}

\begin{proof}
我们沿用引理 6.3.4 的证明中的论证。

注意形如
\begin{equation}
    \prod_{l_a<0}b_{i_a,l_a}\prod_{m_b<0}e^R_{\alpha_b,m_b}\prod_{n_c\le0}a^*_{\alpha_c,n_c}|0\rangle
    \tag{7.1-6}
\end{equation}
的字典序单项式（其中 $e^R_{\alpha,n}$ 由公式 (6.1-3) 给出）构成 $W_{0,\kappa_c}$ 的一组基。于是我们有张量积分解
\[
    W_{0,\kappa_c}=\overline{W}_{0,\kappa_c}\otimes M_{\mathfrak{g},+},
\]
其中 $M_{\mathfrak{g},+}$（相应地，$\overline{W}_{0,\kappa_c}$）是只含 $a^*_{\alpha,n}$（相应地，只含 $e^R_{\alpha,m}$ 与 $b_{i,l}$）的单项式张成的空间。

$\mathfrak{z}(\hat{\mathfrak{g}})$ 的任一元素在 $W_{0,\kappa_c}$ 中的像是 $L_+\mathfrak{g}$-不变向量，这里如前用记号 $L_+\mathfrak{g}=\mathfrak{g}[[t]]$。特别地，它被 $L_+\mathfrak{n}_+$ 与 $\mathfrak{h}$ 零化。由于 $L_+\mathfrak{n}_+$ 与 $e^R_{\alpha,n}$、$b_{i,n}$ 交换，它只沿上述张量积分解的第二个因子 $M_{\mathfrak{g},+}$ 作用。按引理 5.6.4，$M_{\mathfrak{g},+}\simeq\operatorname{Coind}^{L_+\mathfrak{g}}_{L_+\mathfrak{b}_-}\mathbb{C}$ 作为 $L_+\mathfrak{n}_+$-模。因此 $M_{\mathfrak{g},+}$ 中 $L_+\mathfrak{n}_+$-不变元的空间是一维的，由常数张成。

于是 $W_{0,\kappa_c}$ 中 $L_+\mathfrak{n}_+$-不变元的空间等于 $\overline{W}_{0,\kappa_c}$。然而单项式 (7.1-6) 去掉第二个因子后的权等于对应于因子 $e^R_{\alpha,m}$ 的正根之和。这样的单项式是 $\mathfrak{h}$-不变的，当且仅当没有这样的因子出现。因此 $W_{0,\kappa_c}$ 中 $(\mathfrak{h}\oplus L_+\mathfrak{n}_+)$-不变元的子空间是由只含 $b_{i,l}$ 的单项式 (7.1-6) 张成的。这正是子空间 $\pi_0\subset W_{0,\kappa_c}$。
\transnote{卡点：为何 $(\mathfrak{h}\oplus L_+\mathfrak{n}_+)$-不变性迫使 (7.1-6) 中不出现 $e^R_{\alpha,m}$ 因子、只剩 $b_{i,l}$。背景：$e^R_{\alpha,m}$ 带权 $-\alpha$（$\alpha>0$），而 $\mathfrak{h}$ 的权可加，故 (7.1-6) 去掉第二个因子后的权等于各 $e^R_{\alpha,m}$ 所对应的正根之和；$\mathfrak{h}$-不变要求该权为 $0$，只有不出现任何 $e^R$ 因子才可能，剩下的恰是 $\pi_0$ 那一块。指代落点：引理 7.1.2 证明末段，原书 p.202；与引理 6.3.4 的论证平行。[依据：原书 p.202 的 (7.1-6) 与权可加性；§6.1.1 的 $e^R_{\alpha,n}$ 记号][来源：本书 p.202。]}
\end{proof}

至此我们完成了定理 4.3.2 证明的第 1、2 步。

% [Tutor Note: 证明最后一步的关键是：$e^R_{\alpha,m}$ 各带一个正根权，而 $\mathfrak{h}$ 作用的权可加，故 $\mathfrak{h}$-不变迫使所有 $e^R$ 因子都不出现（正根之和不可能为 0）。这与引理 6.3.4 的论证完全平行。核对时注意 (7.1-6) 中三类因子分别对应 $\mathfrak{b}_-$ 部分（$b_{i,l}$）、幂零部分（$e^R$）与对偶 Fock 部分（$a^*$），其张量积分解把 $\pi_0$ 恰好识别为“只含 $b$”的那一块。]

\section{$\mathfrak{sl}_2$ 的情形}

下面执行 §7.1 所述计划的第 3、4 步。这将使我们能把 $\mathfrak{z}(\hat{\mathfrak{g}})$ 在 $w_{\kappa_c}$ 下于 $\pi_0$ 中的像描述为某些算子的核之交。

我们先给出 Wakimoto 模之间交织算子的统一构造，即所谓屏蔽算子（分两类），依据 [Feigin and Frenkel 1990b; Feigin and Frenkel 1990c; Frenkel 1991; Feigin et al. 1994; Frenkel 2005b]。先详细分析 $\mathfrak{g}=\mathfrak{sl}_2$ 的情形。

\begin{keypoints}
    \item \textbf{核心思想：}在 $W_{0,k}$（$\beta\gamma$-系统 $\times$ Heisenberg）上构造与 $\widehat{\mathfrak{sl}}_2$ 作用交换的两类屏蔽算子，并用它们刻画 $V_k(\mathfrak{sl}_2)$。
    \item \textbf{数学工具：}模上的顶点算子 $Y_{V,M}$、OPE 与换位公式 (7.2-2)、Friedan--Martinec--Shenker 玻色化、$\Pi_0$ 中 $a(z)^\gamma$ 的实现。
    \item \textbf{关键结论：}$S_k:W_{0,k}\to W_{-2,k}$（第一类）给出正合列 $0\to V_k(\mathfrak{sl}_2)\to W_{0,k}\to W_{-2,k}\to0$；第二类 $\widetilde{S}_k$ 有同样的核。
    \item \textbf{记号提醒：}下标 $k$ 是水平；$\widetilde{S}_k$ 的临界极限（$k\to-2$，即 $\beta\to0$）取首项得 $\beta\overline{S}$，$\overline{S}$ 才是临界水平可用的算子。
\end{keypoints}

\subsection{与顶点代数模相联系的顶点算子}

需要回忆关于顶点代数模的顶点算子的一些一般结果，依据 [Frenkel et al. 1993], §5.1。设 $V$ 是共形顶点代数（定义见 §3.1.4），$M$ 是 $V$-模，即一个向量空间连同线性映射
\[
    Y_M:V\longrightarrow\operatorname{End}M[[z^{\pm1}]]
\]
满足 [Frenkel and Ben-Zvi 2004] 定义 5.1.1 的公理。特别地，
\[
    Y_M(\omega,z)=\sum_{n\in\mathbb{Z}}L^M_nz^{-n-2}
\]
的 Fourier 系数（其中 $\omega$ 是 $V$ 的共形向量）定义了 Virasoro 代数在 $M$ 上的作用。我们把 $L^M_{-1}$ 记作 $T$。

由公式
\begin{equation}
    Y_{V,M}(A,z)B=e^{zT}Y_M(B,-z)A,\qquad A\in M,\ B\in V
    \tag{7.2-1}
\end{equation}
定义线性映射
\[
    Y_{V,M}:M\longrightarrow\operatorname{Hom}(V,M)[[z^{\pm1}]]
\]
（与命题 2.3.2 的斜对称性质比较）。
\transnote{卡点：(7.2-1) 中的因子 $e^{zT}$ 与右端的宗量 $-z$ 从何而来，为何不是简单的 $Y_M(B,-z)A$。背景：这正是命题 2.3.2 的斜对称性——交换两个顶点算子要付出平移 $e^{zT}$（$T=L^M_{-1}$ 生成沿 $z$ 的平移），并把 $z$ 换成 $-z$，因为 $z^{-n}$ 的展开在两侧重新展开时方向相反；少了 $e^{zT}$，得到的只是形式上的反对称，而非顶点算子意义下的斜对称。指代落点：§7.2.1，(7.2-1)，原书 p.204；命题 2.3.2。(7.2-2)--(7.2-4) 的换位子即由此与 OPE 推出。[依据：原书 p.204 与命题 2.3.2][来源：本书 p.202。]}

这是 [Frenkel et al. 1993] 所引入的交织算子的一个例子。按 [Frenkel et al. 1993] 命题 5.1.2，此映射满足如下性质：对任意 $A\in V$、$B\in M$、$C\in V$，存在元素
\[
    f\in M[[z,w]][z^{-1},w^{-1},(z-w)^{-1}]
\]
使得形式幂级数
\[
    Y_M(A,z)Y_{V,M}(B,w)C,\qquad Y_{V,M}(B,w)Y(A,z)C,
\]
\[
    Y_{V,M}(Y_{V,M}(B,w-z)A,z)C,\qquad Y_{V,M}(Y(A,z-w)B,w)C
\]
分别是 $f$ 在
\[
    M((z))((w)),\qquad M((w))((z)),\qquad M((z))((z-w)),\qquad M((w))((z-w))
\]
中的展开（与 [Frenkel and Ben-Zvi 2004] 推论 3.2.3 比较）。滥用记号，我们将写
\[
    Y_M(A,z)Y_{V,M}(B,w)=Y_{V,M}(Y(A,z-w)B,w),
\]
并像 $M=V$、$Y_{V,M}=Y$ 的情形一样（见 §2.3.3），称此公式为算子乘积展开（OPE）。

在下文中，我们对 $Y(A,z)$ 与 $Y_M(A,z)$ 使用同一记号 $Y(A,z)$。下列 $Y(A,z)$ 与 $Y_{V,M}(B,w)$ 的 Fourier 系数之间的对易关系，其证明与 $M=V$ 的情形完全相同（见 [Frenkel and Ben-Zvi 2004], §3.3.6）。若写
\[
    Y(A,z)=\sum_{n\in\mathbb{Z}}A_{(n)}z^{-n-1},\qquad
    Y_{V,M}(B,w)=\sum_{n\in\mathbb{Z}}B_{(n)}w^{-n-1},
\]
则
\begin{equation}
    [B_{(m)},A_{(k)}]=\sum_{n\ge0}\binom{m}{n}(B_{(n)}A)_{(m+k-n)}.
    \tag{7.2-2}
\end{equation}
特别地，我们得到
\begin{equation}
    \left[\int Y_{V,M}(B,z)dz,\,Y(A,w)\right]
    =Y_{V,M}\left(\int Y_{V,M}(B,z)dz\cdot A,\,w\right).
    \tag{7.2-3}
\end{equation}
这里 $\int$ 如前表示在 $z=0$ 处的留数。

我们需要的另一条性质是公式 (2.3-5) 的如下类比：
\begin{equation}
    Y_{V,M}(TA,z)=\partial_zY_{V,M}(A,z).
    \tag{7.2-4}
\end{equation}
其证明如下：
\begin{align*}
    Y_{V,M}(TA,z)B&=e^{zT}Y_M(B,-z)TA=e^{zT}TY_M(B,-z)A-e^{zT}[T,Y_M(B,-z)]A\\
    &=\partial_z(Y_{V,M}(A,z)B),
\end{align*}
其中用到恒等式
\[
    [T,Y_M(B,z)]=\partial_zY_M(B,z),
\]
它来自 [Frenkel and Ben-Zvi 2004] 命题 5.1.2。

% [Tutor Note: (7.2-1) 中的 $e^{zT}$ 与宗量 $-z$ 是斜对称性的具象化：模上的顶点算子不能直接交换，必须先在 $M$ 一侧作用再移回 $V$ 一侧。核对 (7.2-2) 时请注意它正是 (7.2-1) 与 OPE 的标准推论——换位子沿 $B_{(n)}A$ 的“作用展开”求和，组合系数 $\binom{m}{n}$ 来自 $(z-w)^m$ 的展开方向。]

\subsection{屏蔽算子}

我们把 §7.2.1 的结果用于顶点代数 $W_{0,k}$ 及其模 $W_{-2,k}$（$k\ne-2$）的情形。如前，用 $|0\rangle$ 与 $|-2\rangle$ 记这两个模的最高权向量。

回忆 §6.2.2：在同态 $w_k$ 下，$\widehat{\mathfrak{sl}}_2$ 的生成元 $e_n$ 映到 $a_n$（从现在起，滥用记号，我们把 $\widehat{\mathfrak{sl}}_2$ 的元素与其在 $w_k$ 下的像等同）。$\widehat{\mathfrak{sl}}_2$ 的对易关系给出如下公式：
\[
    [e_n,a_{-1}]=0,\qquad[h_n,a_{-1}]=2a_{n-1},\qquad[f_n,a_{-1}]=-h_{n-1}+k\delta_{n,1}.
\]
此外，由公式 (6.2-3) 可得
\[
    e_n|-2\rangle=a_n|-2\rangle=0,\quad n\ge0;\qquad
    h_n|-2\rangle=f_n|-2\rangle=0,\quad n>0;\qquad
    h_0|-2\rangle=-2|-2\rangle.
\]
因此
\[
    e_n\cdot a_{-1}|-2\rangle=h_n\cdot a_{-1}|-2\rangle=0,\qquad n\ge0,
\]
且
\[
    f_n\cdot a_{-1}|-2\rangle=0,\qquad n>1.
\]
我们还得到
\[
    f_1\cdot a_{-1}|-2\rangle=(k+2)|-2\rangle,
\]
以及
\begin{equation}
    f_0\cdot a_{-1}|-2\rangle=a_{-1}f_0|-2\rangle-h_0|-2\rangle=-b_{-1}|-2\rangle=(k+2)T|-2\rangle,
    \tag{7.2-5}
\end{equation}
其中 $T$ 是平移算子。最后一个等式来自命题 6.2.2 中 $W_{0,k}$ 的共形向量公式，它说明 $T$ 在 $\pi_0\subset W_{0,k}$ 上的作用为
\[
    T=\frac{1}{2(k+2)}\sum_{n\in\mathbb{Z}}b_nb_{-n-1}.
\]

我们希望写出算子
\[
    S_k(z)\overset{\mathrm{def}}{=}Y_{W_{0,k},W_{-2,k}}(a_{-1}|-2\rangle):W_{0,k}\longrightarrow W_{-2,k}
\]
的显式公式，并利用上述性质证明其余数是交织算子，即它与 $\widehat{\mathfrak{sl}}_{2,k}$ 的作用交换。由于 $W_{0,k}$ 的顶点子代数 $M_{\mathfrak{sl}_2}$ 与 $\pi_0^{k+2}$ 彼此交换，我们得到
\[
    S_k(z)=Y_{M_{\mathfrak{sl}_2}}(a_{-1}|0\rangle,z)Y_{\pi_0^{k+2},\pi_{-2}^{k+2}}(|-2\rangle,z)=a(z)Y_{\pi_0^{k+2},\pi_{-2}^{k+2}}(|-2\rangle,z).
\]
还需确定
\[
    V_{-2}(z)\overset{\mathrm{def}}{=}Y_{\pi_0^{k+2},\pi_{-2}^{k+2}}(|-2\rangle,z):\pi_0^{k+2}\longrightarrow\pi_{-2}^{k+2}.
\]
恒等式 (7.2-2) 用于 $A=b_{-1}|0\rangle$ 与 $B=|-2\rangle$ 的情形，给出如下对易关系：
\begin{equation}
    [b_n,V_{-2}(z)]=-2z^nV_{-2}(z).
    \tag{7.2-6}
\end{equation}
此外，由公式 (7.2-4) 与 (7.2-5) 得到
\begin{equation}
    (k+2)\partial_zV_{-2}(z)=-:b(z)V_{-2}(z):.
    \tag{7.2-7}
\end{equation}
容易看出公式 (7.2-6) 与 (7.2-7) 唯一确定 $V_{-2}(z)$（详见 [Frenkel and Ben-Zvi 2004], §5.2.6）。结果得到如下显式公式：
\begin{equation}
    V_{-2}(z)=T_{-2}\exp\left(\frac{1}{k+2}\sum_{n<0}\frac{b_n}{n}z^{-n}\right)
    \exp\left(\frac{1}{k+2}\sum_{n>0}\frac{b_n}{n}z^{-n}\right).
    \tag{7.2-8}
\end{equation}
这里 $T_{-2}$ 表示把最高权向量 $|0\rangle$ 送到最高权向量 $|-2\rangle$、且与所有 $b_n\,(n\ne0)$ 交换的平移算子 $\pi_0^{k+2}\to\pi_{-2}^{k+2}$。
\transnote{卡点：(7.2-7) 的符号、(7.2-8) 指数无负号、(7.2-11) 指数带负号，三处看似矛盾。背景：对 (7.2-8) 直接求导，$\partial_zz^{-n}=-nz^{-n-1}$ 使两个指数因子各贡献 $-\frac{1}{k+2}b_nz^{-n-1}$，故必有 $(k+2)\partial_zV_{-2}(z)=-{:}b(z)V_{-2}(z){:}$，即 (7.2-7) 带负号（原书 p.205 印作带负号；旧译此处脱漏负号，已据页图订正）。(7.2-8) 与 (7.2-11) 的差别则来自字典 $V_{-2}=e^{-(k+2)^{-1}\phi}$ 与 $V_{2(k+2)}=e^{\phi}$：$\phi$ 的系数不同，指数的正负随之不同，并非笔误。指代落点：§7.2.2 (7.2-7)--(7.2-8)，原书 p.205；§7.2.4 (7.2-11)，原书 p.209。[来源：A 级 [Frenkel and Ben-Zvi 2004] §5.2.6（本书正文所引自由场公式）；原书 p.205、p.209]}

现在用公式 (7.2-2) 得到算子 $S_k(z)$ 的如下 OPE：
\[
    e(z)S_k(w)=\mathrm{reg.},\qquad h(z)S_k(w)=\mathrm{reg.},
\]
\[
    f(z)S_k(w)=\frac{(k+2)V_{-2}(w)}{(z-w)^2}+\frac{(k+2)\partial_wV_{-2}(w)}{z-w}+\mathrm{reg.}
    =(k+2)\partial_w\frac{V_{-2}(w)}{z-w}+\mathrm{reg.}
\]
由于全导数的留数为 $0$，由此得到：

\begin{proposition}
    留数 $S_k=\int S_k(w)dw$ 是 $\widehat{\mathfrak{sl}}_2$-模 $W_{0,k}$ 与 $W_{-2,k}$ 之间的交织算子。
\end{proposition}

我们称 $S_k$ 为 $\widehat{\mathfrak{sl}}_2$ 的\textbf{第一类屏蔽算子}。

% [Tutor Note: (7.2-6) 与 (7.2-7) 一起唯一确定 $V_{-2}(z)$：前者对所有 $n$ 给出 $V_{-2}$ 的 $b$-荷为 $-2$，后者把 $T$ 作用与正规序乘积联系起来。核对 (7.2-8) 时注意指数中**没有**负号（与后面的 (7.2-11) 相反），且两个指数分别对 $n<0$ 与 $n>0$ 求和：负号约定一旦写错，$[b_n,V_{-2}]$ 的权就变成 $+2$，整个正合列随之失效。]

\begin{proposition}
    对 $k\notin-2+\mathbb{Q}_{\ge0}$，序列
    \[
        0\longrightarrow V_k(\mathfrak{sl}_2)\longrightarrow W_{0,k}\xrightarrow{\ S_k\ }W_{-2,k}\longrightarrow 0
    \]
    是正合的。
\transnote{卡点：命题 7.2.2 的正合性靠什么成立，以及证明中 $S_k=\sum_na_nV_{-2,-n}$ 的下标约定。背景：三件事各司其职——(i) $S_k$ 是交织算子，故 $\operatorname{Ker}S_k$ 是 $\widehat{\mathfrak{sl}}_2$-子模且含 $V_k(\mathfrak{sl}_2)$；(ii) 用 Kac--Kazhdan 的特征标公式比较 $\operatorname{ch}V_k$ 与 $\operatorname{ch}W_{\lambda,k}$（含因子 $u^{\lambda/2}$），得 $\operatorname{ch}W_{0,k}=\operatorname{ch}V_k+\operatorname{ch}W_{-2,k}$；(iii) $S_k(a_0^*|0\rangle)=|-2\rangle$ 说明 $S_k$ 满射。三者合起来给出 $\operatorname{Ker}S_k=V_k(\mathfrak{sl}_2)$。下标约定：原文以 $V_{-2,n}$ 记 $z^{-n}$ 的系数，译文改记 $V_{-2,-n}$ 为 $z^n$ 的系数，两套编号各自自洽，勿混用。指代落点：命题 7.2.2 及其证明，原书 p.206--207。[来源：A 级 Kac and Kazhdan 1979（本书所引特征标公式）；原书 p.206--207]}
\end{proposition}

\begin{proof}
由命题 7.1.1，对任意 $k$ 值，$V_k(\mathfrak{sl}_2)$ 自然地是 $W_{0,k}$ 的 $\widehat{\mathfrak{sl}}_2$-子模。模 $V_k(\mathfrak{sl}_2)$ 由真空向量 $|0\rangle$ 生成，而它在 $W_{0,k}$ 中的像是最高权向量。我们有
\[
    S_k=\sum_{n\in\mathbb{Z}}a_nV_{-2,-n},
\]
其中 $V_{-2,-n}$ 是 $V_{-2}(z)$ 中 $z^n$ 的系数。由公式 (7.2-8) 显然可知，对所有 $m>0$ 有 $V_{-2,m}|0\rangle=0$。同时对所有 $n\ge0$ 有 $a_n|0\rangle=0$。因此 $|0\rangle$ 属于 $S_k$ 的核。但 $S_k$ 与 $\widehat{\mathfrak{sl}}_2$ 的作用交换，故整个子模 $V_k(\mathfrak{sl}_2)$ 都落在 $S_k$ 的核中。

为证明 $V_k(\mathfrak{sl}_2)$ 与 $S_k$ 的核重合，我们比较二者的特征标（特征标的定义见命题 6.3.3 的证明）。我们将使用记号 $q=e^{-\delta}$、$u=e^{\alpha}$。

由于 $V_k(\mathfrak{sl}_2)$ 同构于由 $n<0$ 的 $e_n,h_n,f_n$ 张成的李代数的泛包络代数，我们得到
\begin{equation}
    \operatorname{ch}V_k(\mathfrak{g})=\prod_{n>0}(1-q^n)^{-1}(1-uq^n)^{-1}(1-u^{-1}q^n)^{-1}.
    \tag{7.2-9}
\end{equation}
类似地得到
\[
    \operatorname{ch}W_{\lambda,k}=u^{\lambda/2}\prod_{n>0}(1-q^n)^{-1}(1-uq^n)^{-1}(1-u^{-1}q^{n-1})^{-1}.
\]
于是 $\operatorname{ch}W_{\lambda,k}=\operatorname{ch}\mathbb{M}_{\lambda,k}$，其中 $\mathbb{M}_{\lambda,k}$ 是 $\widehat{\mathfrak{sl}}_2$ 上最高权为 $(\lambda,k)$ 的 Verma 模。因此若 $\mathbb{M}_{\lambda,k}$ 不可约，则 $W_{\lambda,k}$ 也不可约。使 $\mathbb{M}_{\lambda,k}$ 不可约的 $(\lambda,k)$ 之集在 [Kac and Kazhdan 1979] 中有描述。由该描述可知，若 $k\notin-2+\mathbb{Q}_{\ge0}$，则 $\mathbb{M}_{-2,k}$ 从而 $W_{-2,k}$ 不可约。容易验证 $S_k(a_0^*|0\rangle)=|-2\rangle$，故 $S_k$ 是非零同态。于是对这样的 $k$ 它是满射。因此其核的特征标对这样的 $k$ 等于 $\operatorname{ch}W_{0,k}-\operatorname{ch}W_{-2,k}=\operatorname{ch}V_k(\mathfrak{sl}_2)$。证明完成。
\end{proof}

接下来我们描述 $\widehat{\mathfrak{sl}}_2$ 的第二类屏蔽算子。为此需要回忆 Friedan--Martinec--Shenker 玻色化。

% [Tutor Note: 正合列的存在性靠三件事：(i) $\ker S_k\supset V_k$；(ii) 满射性（$S_k(a_0^*|0\rangle)=|-2\rangle\ne0$）；(iii) 特征标相等。注意 (ii) 用到的 $a_0^*$ 是 βγ 系统的生成元，$a_0^*|0\rangle$ 落在 $W_{0,k}$ 里而不是 $V_k$ 里，这一点是“非平凡满射”的关键。另 (7.2-9) 里原文写 $\operatorname{ch}V_k(\mathfrak{g})$，虽在 $\mathfrak{sl}_2$ 语境下仍用 $\mathfrak{g}$，译文照原文。]

\subsection{Friedan--Martinec--Shenker 玻色化}

考虑以 $p_n,q_n\,(n\in\mathbb{Z})$ 为生成元、以 $\mathbf{1}$ 为中心元的 Heisenberg 李代数，其（非零）对易关系为
\[
    [p_n,p_m]=n\delta_{n,-m}\mathbf{1},\qquad
    [q_n,q_m]=-n\delta_{n,-m}\mathbf{1},\qquad
    [p_n,q_m]=0.
\]
我们令
\[
    p(z)=\sum_{n\in\mathbb{Z}}p_nz^{-n-1},\qquad
    q(z)=\sum_{n\in\mathbb{Z}}q_nz^{-n-1}.
\]
对 $\lambda\in\mathbb{C}$、$\mu\in\mathbb{C}$，设 $\Pi_{\lambda,\mu}$ 是此李代数的 Fock 表示，它由最高权向量 $|\lambda,\mu\rangle$ 生成，满足
\[
    p_n|\lambda,\mu\rangle=\lambda\delta_{n,0}|\lambda,\mu\rangle,\qquad
    q_n|\lambda,\mu\rangle=\mu\delta_{n,0}|\lambda,\mu\rangle,\qquad n\ge0;\qquad
    \mathbf{1}|\lambda,\mu\rangle=|\lambda,\mu\rangle.
\]
考虑顶点算子 $V_{\lambda,\mu}(z):\Pi_{\lambda',\mu'}\to\Pi_{\lambda+\lambda',\mu+\mu'}$，其公式为
\[
    V_{\lambda,\mu}(z)=T_{\lambda,\mu}z^{\lambda\lambda'-\mu\mu'}
    \exp\left(-\sum_{n<0}\frac{\lambda p_n+\mu q_n}{n}z^{-n}\right)
    \exp\left(-\sum_{n>0}\frac{\lambda p_n+\mu q_n}{n}z^{-n}\right),
\]
其中 $T_{\lambda,\mu}$ 是把最高权向量送到最高权向量、且与所有 $p_n,q_n\,(n\ne0)$ 交换的平移算子 $\Pi_{0,0}\to\Pi_{\lambda,\mu}$。

滥用记号，我们将把这些算子写作 $e^{\lambda u+\mu v}$，其中 $u(z)$ 与 $v(z)$ 分别表示 $p(z)$ 与 $q(z)$ 的原函数，即 $p(z)=\partial_zu(z)$、$q(z)=\partial_zv(z)$。

对 $\gamma\in\mathbb{C}$，令
\[
    \Pi_\gamma=\bigoplus_{n\in\mathbb{Z}}\Pi_{n+\gamma,n+\gamma}.
\]
利用这些顶点算子，可按 [Frenkel and Ben-Zvi 2004], §5.2.6 的方式在直和 $\Pi_0$（其真空向量为 $|0,0\rangle$）上定义顶点代数结构。此外，对任意 $\gamma\in\mathbb{C}$，$\Pi_\gamma$ 是顶点代数 $\Pi_0$ 上的模。

下述用顶点代数 $\Pi_0$ 实现顶点代数 $M$（亦称“$\beta\gamma$-系统”）的结果属于 Friedan、Martinec 与 Shenker [1986]。（定理中所述“$M$ 在 $\Pi_0$ 中的像等于 $\int e^udz$ 的核”这一同构，是 [Feigin and Frenkel 1991] 建立的。）

\begin{theorem}
    存在顶点代数的（唯一）嵌入 $M\hookrightarrow\Pi_0$，在该嵌入下，场 $a(z)$ 与 $a^*(z)$ 分别映到场
    \[
        \widetilde{a}(z)=e^{u+v},\qquad
        \widetilde{a}^*(z)=(\partial_ze^{-u})e^{-v}=-:p(z)e^{-u-v}:.
    \]
\end{theorem}

进一步，$M$ 在 $\Pi_0$ 中的像等于算子 $\int e^udz$ 的核。

等价地，$\Pi_0$ 可描述为 $M$ 关于 $a_{-1}$ 的局部化，即
\begin{equation}
    \Pi_0\simeq M[(a_{-1})^{-1}]=\mathbb{C}[a_n]_{n<-1}\otimes\mathbb{C}[a_n^*]_{n\le0}\otimes\mathbb{C}[(a_{-1})^{\pm1}].
    \tag{7.2-10}
\end{equation}
$\Pi_0$ 上的顶点代数结构由 $M$ 上的顶点代数结构自然延拓而得。

在嵌入 $M\hookrightarrow\Pi_0$ 下，§6.2.3 描述的 $M$ 的 Virasoro 场 $T(z)=:\partial_za^*(z)a(z):$ 映到 $\Pi_0$ 中的如下场：
\[
    \frac{1}{2}:p(z)^2:-\frac{1}{2}\partial_zp(z)-\frac{1}{2}:q(z)^2:+\frac{1}{2}\partial_zq(z).
\]
于是映射 $M\hookrightarrow\Pi_0$ 成为共形顶点代数的同态，其中两边的共形结构分别由上述场给出。

FMS 玻色化对我们有用的原因在于：它使我们能赋予场 $a(z)^\gamma$ 意义，其中 $\gamma$ 是任意复数。即，我们把 $a(z)^\gamma$ 换成场
\[
    \widetilde{a}(z)^\gamma=e^{\gamma(u+v)}:\Pi_0\longrightarrow\Pi_\gamma,
\]
它良定义。

现在取张量积 $\Pi_0\otimes\pi_0^{k+2}$，我们再次假设 $k\ne-2$。这是一个顶点代数，含有 $M\otimes\pi_0^{k+2}$，从而含有 $V_k(\mathfrak{sl}_2)$ 作为顶点子代数。特别地，对任意 $\gamma,\lambda\in\mathbb{C}$，张量积 $\Pi_\gamma\otimes\pi_\lambda^{k+2}$ 是 $V_k(\mathfrak{sl}_2)$ 上的模，从而是 $\widehat{\mathfrak{sl}}_2$ 上的模。我们把它记作 $\widetilde{W}_{\gamma,\lambda,k}$。此外，引入玻色顶点算子
\begin{equation}
    V_{2(k+2)}(z)=T_{2(k+2)}
    \exp\left(-\sum_{n<0}\frac{b_n}{n}z^{-n}\right)
    \exp\left(-\sum_{n>0}\frac{b_n}{n}z^{-n}\right).
    \tag{7.2-11}
\end{equation}
令
\begin{equation}
    \widetilde{S}_k(z)=\widetilde{a}(z)^{-(k+2)}V_{2(k+2)}(z).
    \tag{7.2-12}
\end{equation}
与 §7.2.2 中所做的类似计算给出：

\begin{proposition}
    留数
    \[
        \widetilde{S}_k=\int\widetilde{S}_k(z)dz:\widetilde{W}_{0,0,k}\longrightarrow\widetilde{W}_{-(k+2),2(k+2),k}
    \]
    是 $\widehat{\mathfrak{sl}}_2$-模 $\widetilde{W}_{0,0,k}$ 与 $\widetilde{W}_{-(k+2),2(k+2),k}$ 之间的交织算子。
\transnote{卡点：为什么必须先做 FMS 玻色化，才写得出 $a(z)^{\gamma}$。背景：§7.2.2 中 $a(z)$ 只是 Heisenberg 场，其复幂没有先验意义；FMS 实现把 $\Pi_0$ 与 $M[(a_{-1})^{-1}]$ 等同（(7.2-10)），在其中有 $\widetilde{a}(z)=e^{u+v}$，于是 $a(z)^{\gamma}\mapsto\widetilde{a}(z)^{\gamma}=e^{\gamma(u+v)}$ 成为良定义的顶点算子，第二类屏蔽流 $\widetilde{S}_k=\widetilde{a}(z)^{-(k+2)}V_{2(k+2)}(z)$ 才写得出来。这也解释了第二类为何用 $\widetilde{a}$ 的负幂而非 $a$ 的负幂。指代落点：定理 7.2.3（FMS 嵌入）、(7.2-10)--(7.2-12)，原书 p.208--209；第二类屏蔽算子首见于 Dotsenko 1990，见 §7.2.4。[来源：A 级 Friedan--Martinec--Shenker 1986、Feigin--Frenkel 1991（本书 p.208 所引）；Dotsenko 1990]}
\end{proposition}

% [Tutor Note: 这一步是全章的“技术心脏”：βγ 系统的场 $a(z)$ 本身不能取任意复幂，但通过 FMS 玻色化把它写成 $e^{u+v}$ 后，$a(z)^\gamma=e^{\gamma(u+v)}$ 就有意义（代价是换到 $\Pi_\gamma$，即改变 Fock 真空）。注意 (7.2-11) 的指数带负号，与 (7.2-8) 相反；这不是笔误，而是因为 (7.2-11) 中 $V_{2(k+2)}$ 对应的是 $\Pi$ 一侧的真空移动。]

我们称 $\widetilde{S}_k$，或其到 $W_{0,k}\subset\widetilde{W}_{0,0,k}$ 的限制，为 $\widehat{\mathfrak{sl}}_2$ 的\textbf{第二类屏蔽算子}。这个算子最先由 V. Dotsenko [1990] 引入。

\begin{proposition}
    对一般的 $k$，$\widehat{\mathfrak{sl}}_2$-子模 $V_k(\mathfrak{sl}_2)\subset W_{0,k}$ 等于 $\widetilde{S}_k:W_{0,k}\to\widetilde{W}_{-(k+2),2(k+2),k}$ 的核。
\end{proposition}

\begin{proof}
我们将证明：对一般的 $k$，同态
\[
    \widetilde{S}_k:W_{0,k}\longrightarrow\widetilde{W}_{-(k+2),2(k+2),k}
\]
的核与第一类屏蔽算子 $S_k:W_{0,k}\to W_{-2,k}$ 的核重合。这与命题 7.2.2 合起来即得本命题。

注意算子 $S_k$ 显然可延拓为按同一公式定义的算子 $\widetilde{W}_{0,0,k}\to\widetilde{W}_{0,-2,k}$。$\widetilde{S}_k$（相应地，$S_k$）在 $W_{0,k}$ 中的核等于 $W_{0,k}\subset\widetilde{W}_{0,0,k}$ 与 $\widetilde{S}_k$（相应地，$S_k$）作为从 $\widetilde{W}_{0,0,k}$ 到 $\widetilde{W}_{-(k+2),2(k+2),k}$（相应地，$\widetilde{W}_{0,-2,k}$）的算子的核之交。因此只需证明：对一般的 $k$，$\widetilde{S}_k$ 与 $S_k$ 在 $\widetilde{W}_{0,0,k}$ 上的核相等。

现在设 $\phi(z)$ 是 $b(z)$ 的原函数，即 $b(z)=\partial_z\phi(z)$。滥用记号，我们将写 $V_{-2}(z)=e^{-(k+2)^{-1}\phi}$ 与 $V_{2(k+2)}(z)=e^{\phi}$。于是屏蔽流 $S_k(z)$ 与 $\widetilde{S}_k(z)$ 成为
\[
    S_k(z)=e^{u+v-(k+2)^{-1}\phi(z)},\qquad
    \widetilde{S}_k(z)=e^{-(k+2)u-(k+2)v+\phi}.
\]
考虑更一般的情形：设 $\mathfrak{h}$ 是带非退化内积 $\kappa$ 的阿贝尔李代数。用此内积把 $\mathfrak{h}$ 与 $\mathfrak{h}^*$ 等同。设 $\widehat{\mathfrak{h}}_\kappa$ 是 Heisenberg 李代数，$\pi_\chi^\kappa$（$\lambda\in\mathfrak{h}^*\simeq\mathfrak{h}$）是它的 Fock 表示，按 §6.2.1 定义。则对任意 $\chi\in\mathfrak{h}^*$，存在顶点算子 $V_\chi^\kappa(z):\pi_0^\kappa\to\pi_\chi^\kappa$，由公式
\begin{equation}
    V_\chi^\kappa(z)=T_\chi
    \exp\left(-\sum_{n<0}\frac{\chi_n}{n}z^{-n}\right)
    \exp\left(-\sum_{n>0}\frac{\chi_n}{n}z^{-n}\right)
    \tag{7.2-13}
\end{equation}
给出，其中 $\chi_n=\chi\otimes t^n$ 是 Heisenberg 李代数 $\widehat{\mathfrak{h}}_\kappa$ 中对应于 $\chi\in\mathfrak{h}^*$ 的元素（我们用内积 $\kappa$ 把 $\mathfrak{h}^*$ 与 $\mathfrak{h}$ 等同）。

设 $\kappa(\chi,\chi)\ne0$，用 $\check{\chi}$ 记 $\mathfrak{h}$ 中等于 $-2\chi/\kappa(\chi,\chi)$ 的元素。我们断言：若 $\chi$ 是一般的（即远离 $\mathfrak{h}$ 中可数多个超曲面），则 $\int V_\chi^\kappa(z)dz$ 与 $\int V_{\check{\chi}}^\kappa(z)dz$ 在 $\pi_0^\kappa$ 中的核重合。事实上，我们有 $\widehat{\mathfrak{h}}_\kappa$ 的直和分解：由 $\chi_n\,(n\in\mathbb{Z})$ 生成的 Heisenberg 李子代数 $\widehat{\mathfrak{h}}_\kappa^\chi$，以及 $\chi_n\,(n\in\mathbb{Z})$ 在 $\widehat{\mathfrak{h}}_\kappa$ 中的中心化子，即对应于 $\chi$ 在 $\mathfrak{h}$ 中的正交补的 Heisenberg 李子代数 $\widehat{\mathfrak{h}}_\kappa^{\perp}$（注意按对 $\chi$ 的假设，该正交补不包含 $\chi$）。Fock 表示 $\pi_0^\kappa$ 分解为这两个李子代数的 Fock 表示的张量积。算子 $\int V_\chi^\kappa(z)dz$ 与 $\int V_{\check{\chi}}^\kappa(z)dz$ 都与 $\widehat{\mathfrak{h}}_\kappa^{\perp}$ 交换。因此这两个算子在 $\pi_0^\kappa$ 中的核都等于 $\widehat{\mathfrak{h}}_\kappa^{\perp}$ 的 Fock 表示与相应算子在 $\widehat{\mathfrak{h}}_\kappa^\chi$ 的 Fock 表示上的核的张量积。

换言之，核由 Heisenberg 李代数 $\widehat{\mathfrak{h}}_\kappa^\chi$ 的 Fock 表示上相应的核决定。如 [Frenkel and Ben-Zvi 2004], §15.4.15 所示，对一般的 $\kappa(\chi,\chi)$ 值，$\int V_\chi^\kappa(z)dz$ 与 $\int V_{\check{\chi}}^\kappa(z)dz$ 的这两个核重合。因此 $\int V_\chi^\kappa(z)dz$ 与 $\int V_{\check{\chi}}^\kappa(z)dz$ 的核也一般地重合。

现在把此结果用于我们的情形，它对应于以 $\overline{u},\overline{v},\overline{\phi}$ 为基的三维李代数 $\mathfrak{h}$，以及基元素之间如下非零内积：
\[
    \kappa(\overline{u},\overline{u})=-\kappa(\overline{v},\overline{v})=1,\qquad
    \kappa(\overline{\phi},\overline{\phi})=2(k+2).
\]
我们的屏蔽流 $S_k(z)$ 与 $\widetilde{S}_k(z)$ 分别等于 $V_\chi^\kappa(z)$ 与 $V_{\check{\chi}}^\kappa(z)$，其中
\[
    \chi=\overline{u}+\overline{v}-(k+2)^{-1}\overline{\phi},\qquad
    \check{\chi}=-(k+2)\chi=-(k+2)\overline{u}-(k+2)\overline{v}+\overline{\phi}.
\]
因此对一般的 $k$，屏蔽算子 $S_k$ 与 $\widetilde{S}_k$ 的核重合。
\transnote{卡点：$\check{\chi}=-2\chi/\kappa(\chi,\chi)$ 何以可取，所谓“一般的 $k$”回避的是什么。背景：该定义要求 $\kappa(\chi,\chi)\ne0$；使 $\kappa(\chi,\chi)=0$ 的 $\chi$ 落在 $\mathfrak{h}$ 中由内积给出的二次锥面上，取不在此锥面上的 $\chi$ 即可——这就是“一般位置”。对 $\mathfrak{sl}_2$ 取基 $\overline{u},\overline{v},\overline{\phi}$ 使 $\kappa(\overline{u},\overline{u})=-\kappa(\overline{v},\overline{v})=1$、$\kappa(\overline{\phi},\overline{\phi})=2(k+2)$，则 $\chi=\overline{u}+\overline{v}-(k+2)^{-1}\overline{\phi}$ 给出 $\check{\chi}=-(k+2)\chi$，故 $S_k$ 与 $\widetilde{S}_k$ 的核重合。指代落点：§7.2.4 证明末段，原书 p.209--211；同一“一般位置”手法见 [Frenkel and Ben-Zvi 2004] §15.4.15。[来源：A 级 [Frenkel and Ben-Zvi 2004] §15.4.15；原书 p.209--211]}
\end{proof}

% [Tutor Note: 这段论证是“用一般位置的超平面回避”的标准技巧：$\check\chi=-2\chi/\kappa(\chi,\chi)$ 正是使得 $V_{\check\chi}$ 与 $V_\chi$ 在 $\widehat{\mathfrak{h}}_\kappa^\chi$ 的 Fock 表示上有相同核的“反射”向量。原文引用 [Frenkel and Ben-Zvi 2004] §15.4.15 作为这一步的依据。另注意 (7.2-13) 中原文写作“$\pi_\chi^\kappa$，$\lambda\in\mathfrak{h}^*\simeq\mathfrak{h}$”，下标 $\chi$ 与 $\lambda$ 不一致，按上下文应统一为 $\chi$，译文照原文保留。]

\section{任意 $\mathfrak{g}$ 的屏蔽算子}

本节对任意单李代数 $\mathfrak{g}$ 构造 $\hat{\mathfrak{g}}_\kappa$ 的 Wakimoto 模之间的屏蔽算子，并用它们刻画 $W_{0,\kappa}$ 中的 $V_\kappa(\mathfrak{g})$。

\begin{keypoints}
    \item \textbf{核心思想：}把 $\mathfrak{g}$ 沿第 $i$ 个 $\mathfrak{sl}_2$ 子代数做半无限抛物诱导，把 $\mathfrak{sl}_2$ 的两类屏蔽算子搬成 $\hat{\mathfrak{g}}_\kappa$ 的屏蔽算子。
    \item \textbf{数学工具：}抛物子代数 $\mathfrak{p}^{(i)}$ 与其 Levi 子代数、$e_i^R(z)$ 的坐标系实现、FMS 玻色化、$\kappa\to\kappa_c$（即 $\beta\to0$）的极限。
    \item \textbf{关键结论：}第一类 $\overline{S}_{i,\kappa}$ 与第二类 $\widetilde{S}_{i,\kappa}$ 都是交织算子；临界极限 $\overline{S}_i$ 给出 $\mathfrak{z}(\hat{\mathfrak{g}})\subset\bigcap_i\operatorname{Ker}\overline{V}_i[1]$（命题 7.3.6）。
    \item \textbf{记号提醒：}$\beta=\frac12(\kappa-\kappa_c)(h_i,h_i)$；$\check{\alpha}_i=h_i\in\mathfrak{h}$ 是第 $i$ 个余根。
\end{keypoints}

\subsection{抛物诱导}

用 $\mathfrak{sl}_2^{(i)}$ 记 $\mathfrak{g}$ 中由 $e_i,h_i,f_i$ 生成的、同构于 $\mathfrak{sl}_2$ 的李子代数。设 $\mathfrak{p}^{(i)}$ 是由 $\mathfrak{b}_-$ 与 $e_i$ 张成的 $\mathfrak{g}$ 的抛物子代数，$\mathfrak{m}^{(i)}$ 是其 Levi 子代数。于是 $\mathfrak{m}^{(i)}$ 等于 $\mathfrak{sl}_2^{(i)}$ 与 $h_i$ 在 $\mathfrak{h}$ 中的正交补 $\mathfrak{h}_i^\perp$ 的直和。

我们对 $\mathfrak{p}^{(i)}$ 应用 §6.3 中关于半无限抛物诱导的结果。按推论 6.3.2，我们得到从 $\widehat{\mathfrak{sl}}_{2,k}\oplus\widehat{\mathfrak{h}}{}_{i,\kappa_0}^{\perp}$ 的光滑表示范畴（其中 $k$ 与 $\kappa_0$ 满足定理 6.3.1 的条件）到光滑 $\hat{\mathfrak{g}}_{\kappa+\kappa_c}$-模范畴的函子。对 $k$ 与 $\kappa_0$ 的条件是：对应于 $(k+2)$ 与 $\kappa_0$ 的 $\mathfrak{sl}_2^{(i)}$ 上的内积，都是 $\mathfrak{g}$ 上某个不变内积 $(\kappa-\kappa_c)$ 的限制。换言之，$(\kappa-\kappa_c)(h_i,h_i)=2(k+2)$ 且 $\kappa_0=\kappa|_{\mathfrak{h}_i^\perp}$。滥用记号，我们将把 $\kappa_0$ 也写作 $\kappa$。若 $k$ 与 $\kappa$ 满足此条件，则对任意水平为 $k$ 的光滑 $\widehat{\mathfrak{sl}}_2$-模 $R$ 与任意光滑 $\widehat{\mathfrak{h}}{}_\kappa^{\perp}$-模 $L$，张量积 $M_{\mathfrak{g},\mathfrak{p}^{(i)}}\otimes R\otimes L$ 是光滑 $\hat{\mathfrak{g}}_{\kappa+\kappa_c}$-模。

注意上述条件意味着
\begin{equation}
    (\kappa-\kappa_c)(h_i,h_j)=\beta a_{ji},
    \tag{7.3-1}
\end{equation}
其中 $a_{ji}$ 是 $\mathfrak{g}$ 的 Cartan 矩阵的 $(ji)$ 元。

特别地，若取 $R$ 为 $\mathfrak{sl}_2$ 上的 Wakimoto 模 $W_{\lambda,k}$，取 $L$ 为 Fock 表示 $\pi_{\lambda_0}^\kappa$，则相应的 $\hat{\mathfrak{g}}_\kappa$-模同构于 Wakimoto 模 $W_{(\lambda,\lambda_0),\kappa+\kappa_c}$，其中 $(\lambda,\lambda_0)$ 是由 $\lambda$ 与 $\lambda_0$ 作出的 $\mathfrak{g}$ 的权。在此同构下，生成元 $a_{\alpha_i,n}\,(n\in\mathbb{Z})$ 有特殊含义：它们对应 $\hat{\mathfrak{g}}$ 的元素 $e_{i,n}$ 的\emph{右}作用，后者在 §6.1.1 中定义。换言之，把上述模等同起来，迫使我们在 $N_+$ 上选择一套坐标系 $\{y_\alpha\}_{\alpha\in\Delta_+}$，使得 $\rho^R(e_i)=\partial/\partial y_{\alpha_i}$（用 §5.2.5 的记号），从而 $w^R(e_i(z))=a_{\alpha_i}(z)$（用 §6.1.1 的记号）。从现在起，我们把 $w^R(e_i(z))$ 记作 $e_i^R(z)$。对 $N_+$ 上的一般坐标系，有
\begin{equation}
    e_i^R(z)=a_{\alpha_i}(z)+\sum_{\beta\in\Delta_+}P_\beta^{R,i}(a_\alpha^*(z))a_\beta(z)
    \tag{7.3-2}
\end{equation}
（见公式 (6.1-1)）。

现在，$\widehat{\mathfrak{sl}}_2$ 的 Wakimoto 模 $W_{\lambda_1,k}$ 与 $W_{\lambda_2,k}$ 之间的任一交织算子，都给出 $\hat{\mathfrak{g}}_{\kappa+\kappa_c}$ 的 Wakimoto 模 $W_{(\lambda_1,\lambda_0),\kappa+\kappa_c}$ 与 $W_{(\lambda_2,\lambda_0),\kappa+\kappa_c}$ 之间的交织算子，对 $\mathfrak{h}_i^\perp$ 的任意权 $\lambda_0$ 都成立。我们将利用这一事实，以及上一节引入的 $\widehat{\mathfrak{sl}}_2$ 屏蔽算子，来构造 $\hat{\mathfrak{g}}$ 的 Wakimoto 模之间的交织算子。
\transnote{卡点：(7.3-1) 右端是 $\beta a_{ji}$ 而不是 $\beta a_{ij}$，容易被当成下标笔误。背景：$a_{ji}$ 是 Cartan 矩阵的第 $(ji)$ 个元素，方向来自 §6.2.1 的对易关系 $[b_{i,n},b_{j,m}]=n\,a_{ji}\delta_{n,-m}$，即 $h_i$ 与 $e_j$ 的配对；同一方向在 (7.3-14) $\mathbf{D}_{b_{i,m}}\cdot b_{j,n}=a_{ji}\delta_{n,m}$ 中再次出现。对单边（simply-laced）的 $\mathfrak{g}$，Cartan 矩阵对称、$a_{ij}=a_{ji}$，差别不显；一般情形必须照原文写 $a_{ji}$。指代落点：§7.3.1，(7.3-1)；对照 (7.3-14)，原书 p.211 与 p.219。[依据：原书 p.211 (7.3-1)、p.219 (7.3-14)；§6.2.1 (6.2-1)][来源：本书 p.211。]}

% [Tutor Note: (7.3-1) 中 Cartan 矩阵的下标方向极易出错：右端是 $\beta a_{ji}$（$a_{ji}$ 是第 $(ji)$ 个元素，即**转置**后与 $(\kappa-\kappa_c)(h_i,h_j)$ 配对），而 (7.3-14) 中 $\mathbf{D}_{b_{i,m}}\cdot b_{j,n}=a_{ji}\delta_{n,m}$ 用的是同一个方向。若把 $a_{ji}$ 误读成 $a_{ij}$，$\mathfrak{sl}_n$ 之外的 Cartan 矩阵（非对称）立刻出错。]

\subsection{第一类屏蔽算子}

设 $\kappa$ 是 $\mathfrak{g}$ 上的非零不变内积。我们对 $\kappa$ 在 $\mathfrak{h}$ 上的限制使用同一记号。现在，对任意 $\chi\in\mathfrak{h}$，由公式 (7.2-13) 联系一个顶点算子 $V_\chi^\kappa(z):\pi_0^\kappa\to\pi_\chi^\kappa$。对 $\kappa\ne\kappa_c$，令
\[
    S_{i,\kappa}(z)\overset{\mathrm{def}}{=}e_i^R(z)V_{-\alpha_i}^{\kappa-\kappa_c}(z):W_{0,\kappa}\longrightarrow W_{-\alpha_i,\kappa},
\]
其中 $e_i^R(z)$ 由公式 (6.1-1) 给出。注意在 §7.2.1 的记号下，
\begin{equation}
    S_{i,\kappa}(z)=Y_{W_{0,\kappa},W_{-\alpha_i,\kappa}}(e_{i,-1}^R|-\alpha_i\rangle,z).
    \tag{7.3-3}
\end{equation}

按命题 7.2.1 与上述讨论，算子
\begin{equation}
    S_{i,\kappa}=\int S_{i,\kappa}(z)dz:W_{0,\kappa}\longrightarrow W_{-\alpha_i,\kappa}
    \tag{7.3-4}
\end{equation}
是由第 $i$ 个 $\widehat{\mathfrak{sl}}_2$ 子代数的第一类屏蔽算子 $S_k$ 诱导的，其中 $k$ 由公式 $(\kappa-\kappa_c)(h_i,h_i)=2(k+2)$ 确定。于是命题 7.2.1 推出：

\begin{proposition}
    对每个 $i=1,\dots,\ell$，算子 $S_{i,\kappa}$ 都是 $\hat{\mathfrak{g}}_\kappa$-模 $W_{0,\kappa}$ 与 $W_{-\alpha_i,\kappa}$ 之间的交织算子。
\end{proposition}

我们称 $S_{i,\kappa}$ 为 $\hat{\mathfrak{g}}_\kappa$ 的\textbf{第 $i$ 个第一类屏蔽算子}。

回忆由命题 7.1.1，$V_\kappa(\mathfrak{g})$ 自然地是 $W_{0,\kappa}$ 的 $\hat{\mathfrak{g}}_\kappa$-子模与顶点子代数。另一方面，由命题 7.3.1，$S_{i,\kappa}\,(i=1,\dots,\ell)$ 的核之交是 $W_{0,\kappa}$ 的 $\hat{\mathfrak{g}}_\kappa$-子模；又由公式 (7.3-3) 与对易关系 (7.2-3)，它也是 $W_{0,\kappa}$ 的顶点子代数。下述命题在 [Feigin and Frenkel 1999] 中证明（本书不用它）。

\begin{proposition}
    对一般的 $\kappa$，$V_\kappa(\mathfrak{g})$ 等于屏蔽算子 $S_{i,\kappa}\,(i=1,\dots,\ell)$ 的核之交。
\end{proposition}

此外，[Feigin and Frenkel 1999], §3 对一般的 $\kappa$ 构造了 $\hat{\mathfrak{g}}_\kappa$-模的复形 $C_\kappa^\bullet(\mathfrak{g})$。它的第 $i$ 次项是 Wakimoto 模 $W_{w(\rho)-\rho,\kappa}$ 的直和，其中 $w$ 取遍 $\mathfrak{g}$ 的外尔群中长度为 $i$ 的所有元素。其零次上同调同构于 $V_\kappa(\mathfrak{g})$，其余上同调全为零。对 $\mathfrak{g}=\mathfrak{sl}_2$，复形 $C_\kappa^\bullet(\mathfrak{sl}_2)$ 的长度为 $1$，与命题 7.2.2 中出现者相同。一般地，该复形的 $0$ 次项 $C_\kappa^0(\mathfrak{g})$ 是 $W_{0,\kappa}$，$1$ 次项 $C_\kappa^1(\mathfrak{g})$ 是 $\bigoplus_{i=1}^{\ell}W_{-\alpha_i,\kappa}$。零次微分是屏蔽算子 $S_{i,\kappa}$ 之和。

在 [Feigin and Frenkel 1999] 中还解释了如何利用量子群 $U_q(\mathfrak{g})$ 的平凡表示的 Bernstein--Gelfand--Gelfand 分解，把屏蔽算子 $S_{i,\kappa}$ 复合起来构造其他交织算子。粗略地说，屏蔽算子 $S_{i,\kappa}$ 满足 $q$-Serre 关系，即带适当参数 $q$ 的量子化泛包络代数 $U_q(\mathfrak{n}_+)$ 的定义关系。于是对一般的 $\kappa$，我们把 $U_q(\mathfrak{g})$ 上 Verma 模 $M_\lambda$ 中权为 $\mu$ 的奇异向量，对应到交织算子 $W_{\lambda,\kappa}\to W_{\mu,\kappa}$。此算子等于屏蔽流 $S_{i,\kappa}(z)$ 的乘积在构形空间上某个闭链上的积分，该闭链取值于由上述奇异向量自然给出的局部系中。

最简单的算子对应 $U_q(\mathfrak{g})$ 上 Verma 模 $M_0$ 中权为 $-\alpha_i$ 的奇异向量 $f_iv_0$。相应的交织算子不是别的，正是我们的屏蔽算子 (7.3-4)。
\transnote{卡点：命题 7.3.2 中的复形 $\mathcal{C}^\bullet_\kappa(\mathfrak{g})$、BGG 分解与 $q$-Serre 关系只在此处出现一次，后文并不使用，容易误以为它们是第 4 步的工具。背景：这是 Feigin--Frenkel 1999 用半无限 BGG 型复形（各项为 $W_{w(\rho)-\rho,\kappa}$，微分来自 $q$-Serre 关系）构造交织算子的路线；最简单的例子是 $U_q(\mathfrak{g})$ 上 Verma 模 $M_0$ 中权为 $-\alpha_i$ 的奇异向量 $f_iv_0$ 所对应的算子，它正是第一类屏蔽算子 (7.3-4)。本书改用第二类屏蔽算子（其 $\kappa\to\kappa_c$ 极限更易控制），故这段只是文献说明。指代落点：§7.3.2 末，命题 7.3.2，原书 p.213。[来源：A 级 Feigin and Frenkel 1999（本书所引）；原书 p.213]}

% [Tutor Note: 这里的 BGG 分解与 $q$-Serre 关系只在“给出其它交织算子如何构造”的说明中出现，本书不用其证明。注意 $W_{w(\rho)-\rho,\kappa}$ 中 $w(\rho)-\rho$ 是权的平移：对 $\mathfrak{sl}_2$ 只有 $w=1$ 与 $w=s$ 两项，故复形长度为 1，与命题 7.2.2 的正合列一致。]

\subsection{第二类屏蔽算子}

为定义第二类屏蔽算子，需要赋予级数 $(e_i^R(z))^\gamma$ 意义，其中 $\gamma$ 为复数。因此我们在 $N_+$ 上选取一套坐标系，使得 $e_i^R(z)=a_{\alpha_i}(z)$（注意这对所有 $i=1,\dots,\ell$ 不能同时做到）。若我们经由从第 $i$ 个 $\widehat{\mathfrak{sl}}_2$ 子代数的 Wakimoto 模作半无限抛物诱导来定义 $\hat{\mathfrak{g}}$ 上的 Wakimoto 模（见 §7.3.1），则这一点自动成立。

选好这样的坐标系后，我们按 §7.2.3 所述，利用由
\[
    a_{\alpha_i,n},\,a_{\alpha_i,n}^*,\qquad n\in\mathbb{Z},
\]
生成的外尔代数的 Friedan--Martinec--Shenker 玻色化，来定义级数 $a_{\alpha_i}(z)^\gamma$。即，有顶点代数
\[
    \Pi_0^{(i)}=\mathbb{C}[a_{\alpha_i,n}]_{n<-1}\otimes\mathbb{C}[a_{\alpha_i,n}^*]_{n\le0}\otimes\mathbb{C}[a_{\alpha_i,-1}^{\pm1}],
\]
它含有
\[
    M_{\mathfrak{g}}^{(i)}=\mathbb{C}[a_{\alpha_i,n}]_{n\le-1}\otimes\mathbb{C}[a_{\alpha_i,n}^*]_{n\le0},
\]
以及按 §7.2.3 定义的 $\Pi_0^{(i)}$-模 $\Pi_\gamma^{(i)}$。然后令
\[
    \widetilde{W}_{\gamma,\lambda,\kappa}^{(i)}=W_{\lambda,\kappa}\otimes_{M_{\mathfrak{g}}^{(i)}}\Pi_\gamma^{(i)}.
\]
这是 $\hat{\mathfrak{g}}_\kappa$-模，当 $\gamma=0$ 时含有 $W_{\lambda,\kappa}$。注意 $\widetilde{W}_{0,0,\kappa_c}^{(i)}$ 是由 $\widehat{\mathfrak{sl}}_2$-模 $\widetilde{W}_{0,0,-2}$ 经半无限抛物诱导所得的 $\hat{\mathfrak{g}}$-模。

现在令 $\beta=\frac{1}{2}(\kappa-\kappa_c)(h_i,h_i)$，并定义场
\begin{equation}
    \widetilde{S}_{\kappa,i}(z)\overset{\mathrm{def}}{=}(e_i^R(z))^{-\beta}V_{\check{\alpha}_i}(z):
    \widetilde{W}_{0,0,\kappa}^{(i)}\longrightarrow\widetilde{W}_{-\beta,\beta\check{\alpha}_i,\kappa}^{(i)}.
    \tag{7.3-5}
\end{equation}
这里 $\check{\alpha}_i=h_i\in\mathfrak{h}$ 表示 $\mathfrak{g}$ 的第 $i$ 个余根。于是算子
\[
    \widetilde{S}_{\kappa,i}=\int\widetilde{S}_{\kappa,i}(z)dz
\]
由第 $i$ 个 $\widehat{\mathfrak{sl}}_2$ 子代数的第二类屏蔽算子 $\widetilde{S}_k$ 诱导，其中 $k$ 由公式 $(\kappa-\kappa_c)(h_i,h_i)=2(k+2)$ 确定。因此命题 7.2.1 推出（类似结果也在 [Petersen et al. 1997] 中得到）：

\begin{proposition}
    算子 $\widetilde{S}_{\kappa,i}$ 是 $\hat{\mathfrak{g}}_\kappa$-模 $\widetilde{W}_{0,0,\kappa}^{(i)}$ 与 $\widetilde{W}_{-\beta,\beta\check{\alpha}_i,\kappa}^{(i)}$ 之间的交织算子。
\end{proposition}

注意 $\widetilde{S}_{\kappa,i}$ 是 $Y_{V,M}((e_{i,-1}^R)^{-\beta}|\check{\alpha}_i\rangle,z)$ 的留数，其中 $V=\widetilde{W}_{0,0,\kappa}^{(i)}$、$M=\widetilde{W}_{-\beta,\beta\check{\alpha}_i,\kappa}^{(i)}$（见 §7.2.1）。因此，按对易关系 (7.2-3)，$\widetilde{S}_{\kappa,i}\,(i=1,\dots,\ell)$ 的核之交自然地是 $\widetilde{W}_{0,0,\kappa}^{(i)}$ 或 $W_{0,\kappa}$ 的顶点子代数。结合命题 7.3.2 与命题 7.2.5，我们得到：

\begin{proposition}
    对一般的 $\kappa$，$V_\kappa(\mathfrak{g})$ 作为 $\hat{\mathfrak{g}}_\kappa$-模与作为顶点代数，都同构于屏蔽算子
    \[
        \widetilde{S}_{i,\kappa}:W_{0,\kappa}\longrightarrow\widetilde{W}_{-\beta,\beta\check{\alpha}_i,\kappa}^{(i)},\qquad i=1,\dots,\ell
    \]
    的核之交。
\end{proposition}

我们指出，沿用 [Feigin and Frenkel 1999] 的步骤，可以利用屏蔽算子 $\widetilde{S}_{i,\kappa}$ 构造更一般的交织算子。
\transnote{卡点：(7.3-5) 中的 $V_{\check{\alpha}_i}(z)$ 没写 $\kappa$ 上标，而本节其他顶点算子都带（如 (7.3-3) 的 $V^{\kappa-\kappa_c}_{-\alpha_i}$）。背景：这里用的是 (7.2-13) 的 $V^\kappa_\chi(z):\pi^\kappa_0\to\pi^\kappa_\chi$，按上下文应为 $V^\kappa_{\check{\alpha}_i}(z)$，其中余根 $\check{\alpha}_i=h_i$；原书 (7.3-5) 省略上标，属记号不统一而非另有所指。指代落点：§7.3.3，(7.3-5)，原书 p.214；对照 (7.2-13)，原书 p.209。译文照原文保留，读者按此补足即可。[依据：原书 p.214 与 p.209 的上标用法][来源：本书 p.213。]}

% [Tutor Note: (7.3-5) 中 $V_{\check{\alpha}_i}(z)$ 没有写上标 $\kappa$，而第一类屏蔽算子的 (7.3-3) 用的是 $V_{-\alpha_i}^{\kappa-\kappa_c}$。按 §7.3.2 的约定，此处应理解为级为 $\kappa$（更确切地 $\kappa-\kappa_c$）的顶点算子；原文在此处省略上标，属记号不统一。另 $\beta=\frac12(\kappa-\kappa_c)(h_i,h_i)$ 中的 $\frac12$ 与 (7.3-1) 的归一化配套：它使得 $\beta=-\frac12(k+2)$，即 $\kappa\to\kappa_c$ 对应 $\beta\to0$。]

\subsection{临界水平下的第二类屏蔽算子}

我们想用屏蔽算子 $\widetilde{S}_{i,\kappa}$ 来刻画同态 $V_{\kappa_c}(\mathfrak{g})\to W_{0,\kappa_c}$ 的像（即实施 §7.1 计划的第 3 步）。首先需要定义它们在 $\kappa\to\kappa_c$ 时的极限。

从 $\mathfrak{g}=\mathfrak{sl}_2$ 的情形开始。为定义 $\widetilde{S}_k$ 当 $k\to-2$ 时的极限，我们把 $W_{0,k}$ 与 $\widetilde{W}_{-(k+2),2(k+2),k}$ 做成 $\mathbb{C}[\beta]$ 上的自由模，其中 $\beta$ 是表示 $k+2$ 的形式变量，然后取这些模对由 $\beta$ 生成的理想的商。

更确切地说，设 $\pi_0[\beta]$（相应地，$\pi_{2\beta}[\beta]$）是由 $b_n\,(n<0)$ 的单项式作用在向量 $|0\rangle$（相应地，$|2\beta\rangle$）上张成的自由 $\mathbb{C}[\beta]$-模。我们按 §6.2.1 的方式在 $\pi_0[\beta]$ 上定义 $\mathbb{C}[\beta]$ 上的顶点代数结构。则 $\pi_{2\beta}[\beta]$ 是 $\pi_0[\beta]$ 上的模。对 $\beta$ 的依赖来自对易关系
\[
    [b_n,b_m]=2\beta n\delta_{n,-m}
\]
以及 $b_0$ 在 $\pi_{2\beta}[\beta]$ 上按乘以 $2\beta$ 作用这一事实。把 $\pi_0[\beta]$（相应地，$\pi_{2\beta}[\beta]$）模去由 $(\beta-k)$（$k\in\mathbb{C}$）生成的理想，便得到 §6.2.2 引入的顶点代数 $\pi_0^k$（相应地，$\pi_0^k$ 上的模 $\pi_{2k}^k$）。

我们定义自由 $\mathbb{C}[\beta]$-模 $W_0[\beta]$ 与 $\widetilde{W}_{0,0}[\beta]$ 为张量积
\[
    M\otimes_{\mathbb{C}}\pi_0[\beta]\qquad\text{和}\qquad\Pi_0\otimes_{\mathbb{C}}\pi_0[\beta].
\]
它们分别是 $\mathbb{C}[\beta]$ 上的顶点代数，且模去由 $(\beta-k)\,(k\in\mathbb{C})$ 生成的理想后分别是顶点代数 $W_{0,k}$ 与 $\widetilde{W}_{0,0,k}$。

其次，设 $\Pi_{-\beta+n,-\beta+n}$ 是由 $p_n,q_n\,(n<0)$ 的字典序单项式作用在向量 $|-\beta+n,-\beta+n\rangle$ 上张成的自由 $\mathbb{C}[\beta]$-模。令
\[
    \Pi_{-\beta}=\bigoplus_{n\in\mathbb{Z}}\Pi_{-\beta+n,-\beta+n},
\]
以及
\[
    \widetilde{W}_{-\beta,2\beta}=\Pi_{-\beta}\otimes_{\mathbb{C}[\beta]}\pi_{2\beta}[\beta].
\]

形式幂级数
\[
    V_{2\beta}(z):\pi_0[\beta]\to\pi_{2\beta}[\beta]\otimes_{\mathbb{C}}\mathbb{C}[\beta],\qquad
    \widetilde{a}(z)^{-\beta}:\Pi_0\otimes_{\mathbb{C}}\mathbb{C}[\beta]\to\Pi_{-\beta}
\]
（由上列公式给出）的每个 Fourier 系数都是良定义的线性算子，且与 $\mathbb{C}[\beta]$ 的作用交换。因此它们乘积的 Fourier 系数也是从 $\widetilde{W}_0[\beta]$ 到 $\widetilde{W}_{-\beta,2\beta}$ 的良定义线性算子。这些算子在由 $(\beta-k)\,(k\in\mathbb{C})$ 生成的理想的商上的相应算子，与上面引入的算子一致。我们需要显式计算留数 $\int\widetilde{a}(z)^{-\beta}V_{2\beta}(z)dz$ 的 $\beta$-展开的首项及其到 $W_0[\beta]$ 的限制。我们将把这个首项用作临界水平的屏蔽算子。

先考虑 $V_{2\beta}(z)$ 按 $\beta$ 的幂的展开。写 $V_{2\beta}(z)=\sum_{n\in\mathbb{Z}}V_{2\beta}[n]z^{-n}$。通过形式幂级数
\begin{equation}
    \sum_{n\le0}\overline{V}[n]z^{-n}=\exp\left(\sum_{m>0}\frac{b_{-m}}{m}z^m\right)
    \tag{7.3-6}
\end{equation}
引入算子 $\overline{V}[n]\,(n\le0)$。

利用公式 (7.2-11)，得到 $V_{2\beta}[n]$ 的如下展开：
\[
    V_{2\beta}[n]=
    \begin{cases}
        \overline{V}[n]+\beta(\dots), & n\le0,\\
        \beta\overline{V}[n]+\beta^2(\dots), & n>0,
    \end{cases}
\]
其中
\begin{equation}
    \overline{V}[n]=-2\sum_{m\le0}\overline{V}[m]\frac{\partial}{\partial b_{m-n}},\qquad n>0.
    \tag{7.3-7}
\end{equation}

其次考虑 $\widetilde{a}(z)^{-\beta}=e^{-\beta(u+v)}$ 按 $\beta$ 的幂的展开。写 $\widetilde{a}(z)^{-\beta}=\sum_{n\in\mathbb{Z}}\widetilde{a}(z)^{-\beta}_{[n]}z^{-n}$。我们将把 $\Pi_{-\beta}$ 与 $\Pi_0\otimes_{\mathbb{C}}\mathbb{C}[\beta]$ 作为 $\mathbb{C}[\beta]$-模等同。于是得到
\[
    \widetilde{a}(z)^{-\beta}_{[n]}=
    \begin{cases}
        1+\beta(\dots), & n=0,\\
        \beta\dfrac{p_n+q_n}{n}+\beta^2(\dots), & n\ne0.
    \end{cases}
\]
上述公式推出屏蔽算子 $\widetilde{S}_{\beta-2}$ 的如下展开：
\[
    \int\widetilde{a}(z)^{-\beta}V_{2\beta}(z)dz=\beta\left(\overline{V}[1]+\sum_{n>0}\frac{1}{n}\overline{V}[-n+1](p_n+q_n)\right)+\beta^2(\dots).
\]
因此我们把 $\beta=0$（对应 $k=-2$）处屏蔽算子 $\widetilde{S}_{\beta-2}$ 的极限定义为算子
\begin{equation}
    \overline{S}\overset{\mathrm{def}}{=}\overline{V}[1]+\sum_{n>0}\frac{1}{n}\overline{V}[-n+1](p_n+q_n),
    \tag{7.3-8}
\end{equation}
它从 $W_{0,-2}=M_{\mathfrak{sl}_2}\otimes\pi_0$ 作用到 $\widetilde{W}_{0,0,-2}=\Pi_0\otimes\pi_0$。由构造，$\overline{S}$ 等于屏蔽算子 $\widetilde{S}_{\beta-2}$ 的 $\beta$-展开的首项。因此 $\overline{S}$ 与 $W_{0,-2}$ 及 $\widetilde{W}_{0,0,-2}$ 上的 $\widehat{\mathfrak{sl}}_2$ 作用交换。

可以把算子
\[
    \frac{p_n+q_n}{n}=-(u_n+v_n),\qquad n>0,
\]
用由 $a_m,a_m^*\,(m\in\mathbb{Z})$ 生成的 Heisenberg 代数表出。即，由定义 $\widetilde{a}(z)=e^{u+v}$ 可得 $u(z)+v(z)=\log\widetilde{a}(z)$，于是场 $u(z)+v(z)$ 与 $a(w)$ 交换，且与 $a^*(w)$ 有如下 OPE：
\[
    (u(z)+v(z))a^*(w)=\frac{1}{z-w}\widetilde{a}(w)^{-1}+\mathrm{reg.}
\]
因此，若写 $\widetilde{a}(z)^{-1}=\sum_{n\in\mathbb{Z}}\widetilde{a}(z)^{-1}_{[n]}z^{-n}$，便得到如下对易关系：
\begin{equation}
    \left[\frac{p_n+q_n}{n},\,a_m^*\right]=-\widetilde{a}(z)^{-1}_{[n+m-1]},\qquad n>0.
    \tag{7.3-9}
\end{equation}
利用 (7.2-10) 中 $\Pi_0$ 的实现，级数 $\widetilde{a}(z)^{-1}$ 可表为
\begin{equation}
    \widetilde{a}(z)^{-1}=(a_{-1})^{-1}\left(1+(a_{-1})^{-1}\sum_{n\ne-1}a_nz^{-n-1}\right)^{-1},
    \tag{7.3-10}
\end{equation}
其中右端按 $(a_{-1})^{-1}$ 的正幂展开为形式幂级数。容易看出此幂级数的每个 Fourier 系数在 $\Pi_0$ 上都是良定义的线性算子。

% [Tutor Note: (7.3-9) 与 (7.3-10) 是全节最易算错的一处。要点有三：其一，$u(z)+v(z)=\log\widetilde{a}(z)$ 中的对数只在 $\Pi_0$ 的局部化实现 (7.2-10) 下才有意义，这就是 $\widetilde{a}(z)^{-1}=(a_{-1})^{-1}(\cdots)^{-1}$ 里必须出现 $(a_{-1})^{-1}$ 的原因；其二，右端必须按 $(a_{-1})^{-1}$ 的**正**幂展开，才能保证每个 Fourier 系数作用在 $W_{0,-2}$ 的任一向量上只产生有限和（原文在 (7.3-10) 之后特意强调这一点）；其三，(7.3-9) 的指标 $n+m-1$ 中 $-1$ 来自 $u(z)+v(z)$ 的 OPE 单极点系数，写错一处会使 $\overline{S}$ 的常数项错位。]

上述公式完全确定了 $(p_n+q_n)/n\,(n>0)$ 从而 $\overline{S}$ 在 $W_{0,-2}=M\otimes\pi_0\subset\widetilde{W}_{0,0,-2}$ 中任一向量上的作用。即，我们用对易关系 (7.3-9) 把 $(p_n+q_n)/n$ 穿过 $a_m^*$；结果得到 $\widetilde{a}(z)^{-1}$ 的 Fourier 系数，它们由公式 (7.3-10) 给出。把其中每个作用于 $W_{0,-2}$ 中任一向量，总得到有限和。
\transnote{卡点：$u+v=\log\widetilde{a}(z)$ 这类对数式为何只能“形式地”写，(7.3-10) 又为何按 $(a_{-1})^{-1}$ 的正幂展开。背景：对数只有在 $\Pi_0\cong M[(a_{-1})^{-1}]$ 这个局部化实现（(7.2-10)）中才有意义——一般地 $\widetilde{a}(z)$ 不可逆，须先局部化使 $\widetilde{a}(z)^{-1}$ 存在，$\log$ 才可展开。同理 (7.3-10) 把 $\widetilde{a}(z)^{-1}$ 写成 $(a_{-1})^{-1}$ 的幂级数，这是形式展开而非收敛展开；正因为 $(a_{-1})^{-1}$ 的正幂只有有限多个落在任一向量上，才总有“有限和”。指代落点：§7.3.4，(7.3-9)--(7.3-10)，原书 p.216--217；对照 §7.2.3 (7.2-10)，原书 p.208。[依据：原书 p.208 (7.2-10)、p.216--217][来源：本书 p.216。]}

这就完成了 $\mathfrak{sl}_2$ 情形下屏蔽算子 $\widetilde{S}_k$ 当 $k\to-2$ 时的极限 $\overline{S}$ 的构造。现在考虑任意 $\mathfrak{g}$ 的情形。

$\widetilde{S}_{i,\kappa}$ 当 $\kappa\to\kappa_c$ 时的极限按定义是由半无限抛物诱导从 $\overline{S}$ 得到的算子
\[
    \overline{S}_i:W_{0,\kappa_c}\longrightarrow\widetilde{W}_{0,0,\kappa_c}^{(i)},
\]
因此它由公式
\begin{equation}
    \overline{S}_i=\overline{V}_i[1]+\sum_{n>0}\frac{1}{n}\overline{V}_i[-n+1](p_{i,n}+q_{i,n})
    \tag{7.3-11}
\end{equation}
给出。这里 $\overline{V}_i[n]:\pi_0\to\pi_0$ 是由公式
\begin{equation}
    \sum_{n\le0}\overline{V}_i[n]z^{-n}=\exp\left(\sum_{m>0}\frac{b_{i,-m}}{m}z^m\right),
    \tag{7.3-12}
\end{equation}
\begin{equation}
    \overline{V}_i[1]=-\sum_{m\le0}\overline{V}_i[m]\mathbf{D}_{b_{i,m-1}},
    \tag{7.3-13}
\end{equation}
给出的线性算子，其中 $\mathbf{D}_{b_{i,m}}$ 表示沿 $b_{i,m}$ 方向的导数，由公式
\begin{equation}
    \mathbf{D}_{b_{i,m}}\cdot b_{j,n}=a_{ji}\delta_{n,m}
    \tag{7.3-14}
\end{equation}
给出，而 $(a_{kl})$ 是 $\mathfrak{g}$ 的 Cartan 矩阵（它按 $\mathbf{D}_{b_{i,m}}\cdot b_{i,n}=2\delta_{n,m}$ 归一化，与 $\mathfrak{sl}_2$ 情形一致）。这来自 $b_{i,n}$ 之间的对易关系 (6.2-1) 与公式 (7.3-1)。

作用在 $\Pi_0^{(i)}$ 上的算子 $(p_{i,n}+q_{i,n})/n$ 按与上面相同的方式定义。

于是我们得到良定义的线性算子 $\overline{S}_i:W_{0,\kappa_c}\to\widetilde{W}_{0,0,\kappa_c}^{(i)}$。由构造，它们与两个模上的 $\hat{\mathfrak{g}}_{\kappa_c}$ 作用交换。显然，算子 $\overline{S}_i\,(i=1,\dots,\ell)$ 零化 $W_{0,\kappa_c}$ 的最高权向量。因此它们零化由最高权向量在 $\hat{\mathfrak{g}}_{\kappa_c}$ 作用下得到的所有向量，即 $V_{\kappa_c}(\mathfrak{g})\subset W_{0,\kappa_c}$ 中的所有向量。于是得到

\begin{proposition}
    真空模 $V_{\kappa_c}(\mathfrak{g})$ 在 $w_{\kappa_c}$ 下的像包含在算子 $\overline{S}_i:W_{0,\kappa_c}\to\widetilde{W}_{0,0,\kappa_c}^{(i)}\,(i=1,\dots,\ell)$ 的核之交中。
\end{proposition}

由命题 7.3.5、命题 7.1.1 与引理 7.1.2 可知，$\mathfrak{z}(\hat{\mathfrak{g}})$ 在同态 $w_{\kappa_c}:V_{\kappa_c}(\mathfrak{g})\hookrightarrow W_{0,\kappa_c}$ 下的像，包含在算子 $\overline{S}_i\,(i=1,\dots,\ell)$ 限制到 $\pi_0\subset W_{0,\kappa_c}$ 后的核之交中。但按公式 (7.3-11)，$\overline{S}_i$ 到 $\pi_0$ 的限制正是由公式 (7.3-13) 给出的算子 $\overline{V}_i[1]:\pi_0\to\pi_0$。于是得到：

\begin{proposition}
    $V_{\kappa_c}(\mathfrak{g})$ 的中心 $\mathfrak{z}(\hat{\mathfrak{g}})$ 包含在 $\pi_0$ 中算子 $\overline{V}_i[1]\,(i=1,\dots,\ell)$ 的核之交中。
\end{proposition}

这就完成了 §7.1 所述计划的第 4 步。
下一章我们将利用此结果来描述中心 $\mathfrak{z}(\hat{\mathfrak{g}})$ 并把它与 opers 空间上的函数代数等同起来。\transnote{卡点：临界水平取的是 $\beta$-展开的首项 $\overline{S}$，而不是 $\int\widetilde{a}(z)^{-\beta}V_{2\beta}(z)dz$ 这个算子本身。背景：$\beta\to0$ 时整体算子并不收敛，它正比于 $\beta$，只有首项有限；同时 $\mathbb{C}[\beta]$ 上的模须模去 $(\beta)$ 才落回 $\beta=0$ 的 $\pi_0$，故只能逐项取极限。于是 $\overline{S}$（及各 $\overline{V}_i[1]$）才是 $\kappa=\kappa_c$ 时的第二类屏蔽算子，命题 7.3.5、7.3.6 说的都是它。指代落点：§7.3.4 末，(7.3-8)，原书 p.218。[依据：原书 p.218 与本节前述 $\beta$-展开][来源：本书 p.217。]}

% [Tutor Note: 临界极限的机制值得单独记住：$\widetilde{S}_{\beta-2}$ 本身在 $\beta\to0$ 时没有极限（它正比于 $\beta$），真正的算子是其首项 $\overline{S}$。原文第 217 页（PDF p.229）明确说“we will use this leading term as the screening operator at the critical level”。因此 $\overline{V}_i[1]$ 只是 $\overline{S}_i$ 限制到 $\pi_0$ 的那一项，这正是命题 7.3.6 的内容。]

\subsection{屏蔽流的其他 Fourier 分量}

下文 §9.6 将需要上面构造的第二类屏蔽流的其他 Fourier 分量（对第一类屏蔽流可沿同样思路得到类似结果）。

从 $\mathfrak{g}=\mathfrak{sl}_2$ 的情形开始。考虑屏蔽流
\[
    \widetilde{S}_k(z)=\widetilde{a}(z)^{-(k+2)}V_{2(k+2)}(z),
\]
它在公式 (7.2-12) 中引入。在命题 7.2.4 中我们考虑了它作为线性算子 $\widetilde{W}_{0,0,k}\to\widetilde{W}_{-(k+2),2(k+2),k}$ 的留数。特别地，$V_{2(k+2)}(z)$ 在公式 (7.2-11) 中定义为从 $\pi_0^{(k+2)}$ 作用到 $\pi_{2(k+2)}^{k+2}$ 的流。

我们也可以考虑此流在 Heisenberg 代数的其他 Fock 模 $\pi_\lambda^{k+2}\,(\lambda\in\mathbb{C})$ 上的作用。不过，此时还须把公式 (7.2-11) 给出的流乘以因子 $z^{b_0}=z^\lambda$。只有这样，上面得到的此流与 $b(z)$ 之间的关系才成立。（若 $\lambda=0$，则 $z^\lambda=1$，这就是为什么我们在公式 (7.2-11) 中没有把它写进去。）从现在起，我们把公式 (7.2-11) 乘以 $z^{b_0}=z^\lambda$ 后的流记作 $V_{2(k+2)}(z)$。它是 $z^\alpha$（其中 $\alpha\in\lambda+\mathbb{Z}$）的形式线性组合。因此对非整数的 $\lambda$ 我们不能取它的留数。但若 $\lambda\in\mathbb{Z}$，则得到良定义的线性算子
\begin{equation}
    \widetilde{S}^{(\lambda)}_k\overset{\mathrm{def}}{=}\int\widetilde{S}_k(z)dz:
    \widetilde{W}_{\lambda,k}\longrightarrow\widetilde{W}_{-(k+2),\lambda+2(k+2),k},
    \tag{7.3-15}
\end{equation}
它同样与 $\widehat{\mathfrak{sl}}_2$ 的作用交换。

其次考虑 $k=\beta-2$ 且 $\beta\to0$ 时此算子的极限。设 $\lambda\in\mathbb{Z}_+$。则与 §7.3.4 中完全一样，可得 $\widetilde{S}^{(\lambda)}_{\beta-2}\to\beta\overline{S}^{(\lambda)}$，其中
\begin{equation}
    \overline{S}^{(\lambda)}=\overline{V}[\lambda+1]+\sum_{n>0}\frac{1}{n}\overline{V}[-n+1](p_{n+\lambda}+q_{n+\lambda}),
    \tag{7.3-16}
\end{equation}
它从 $W_{\lambda,-2}=M_{\mathfrak{sl}_2}\otimes\pi_\lambda$ 作用到 $\widetilde{W}_{0,\lambda,-2}=\Pi_0\otimes\pi_\lambda$。作用在 $\pi_\lambda$ 上的算子 $\overline{V}[n]$ 由公式 (7.3-6) 与 (7.3-7) 定义，而从 $M_{\mathfrak{sl}_2}$ 作用到 $\Pi_0$ 的算子 $p_n+q_n$ 由公式 (7.3-9) 与 (7.3-10) 定义。

由于 $\overline{S}^{(\lambda)}$ 是作为交织算子的极限得到的，我们得到 $\overline{S}^{(\lambda)}$ 也是交织算子（即它与 $\widehat{\mathfrak{sl}}_2$ 的作用交换）。

现在，对任意 $\mathfrak{g}$，我们按 §7.3.4 的方式，把半无限抛物诱导作用于算子 (7.3-15) 与 (7.3-16)，以构造交织算子。首先，对任意满足 $\lambda_i=\langle\lambda,\check{\alpha}_i\rangle\in\mathbb{Z}$ 的权 $\lambda\in\mathfrak{h}^*$，得到交织算子
\[
    \widetilde{S}^{(\lambda_i)}_{\kappa,i}=\int\widetilde{S}_{\kappa,i}(z)dz:
    \widetilde{W}_{\lambda,\kappa}^{(i)}\longrightarrow\widetilde{W}_{-\beta,\lambda+\beta\check{\alpha}_i,\kappa}^{(i)},
\]
其中 $\widetilde{S}_{\kappa,i}(z)$ 是公式 (7.3-5) 给出的流乘以因子 $z^{\lambda_i}$。

其次，考虑此算子当 $\kappa\to\kappa_c$、$\lambda_i\in\mathbb{Z}_+$ 时的极限。如前，这对应 $\beta\to0$，其中 $\beta=\frac{1}{2}(\kappa-\kappa_c)(h_i,h_i)$。极限算子由公式
\begin{equation}
    \overline{S}^{(\lambda_i)}_i=\overline{V}_i[\lambda_i+1]+\sum_{n>0}\frac{1}{n}\overline{V}_i[-n+1](p_{i,n+\lambda_i}+q_{i,n+\lambda_i})
    \tag{7.3-17}
\end{equation}
给出，其中 $\overline{V}_i[n]\,(n\le0)$ 由公式 (7.3-12) 给出，且
\begin{equation}
    \overline{V}_i[\lambda_i+1]=-\sum_{m\le\lambda_i}\overline{V}_i[m-\lambda_i]\mathbf{D}_{b_{i,m-1}},
    \tag{7.3-18}
\end{equation}
其中如前，$\mathbf{D}_{b_{i,m}}$ 是 $\pi_\lambda\simeq\mathbb{C}[b_{i,n}]_{i=1,\dots,\ell;\,n<0}$ 上的导子，满足
\[
    \mathbf{D}_{b_{i,m}}\cdot b_{j,n}=a_{ji}\delta_{n,m}.
\]
算子 $\overline{S}^{(\lambda_i)}_i$ 从 $W_{\lambda,\kappa_c}$ 作用到 $\widetilde{W}_{0,\lambda,\kappa_c}$。特别地，我们发现 $\overline{S}_i$ 到 $\pi_\lambda\subset W_{\lambda,\kappa_c}$ 的限制等于公式 (7.3-18) 给出的 $\overline{V}_i[\lambda_i+1]$。

我们将把此结果用于下文命题 9.6.7 的证明，该命题把命题 7.3.6 推广到任意整优权。
\transnote{卡点：(7.3-16)、(7.3-17) 里 $p_{n+\lambda}$、$p_{i,n+\lambda_i}$ 的指标为何整体平移 $\lambda$。背景：$\widetilde{S}_{\kappa,i}(z)$ 乘上因子 $z^{\lambda_i}$ 后，取 $dz$ 的留数只有 $z^{-1}$ 项存活；$z^{\lambda_i}$ 与各场的 $z^{-n}$ 展开相乘使所有指标平移 $\lambda_i$，于是出现 $\overline{V}_i[\lambda_i+1]$ 与 $(p_{i,n+\lambda_i}+q_{i,n+\lambda_i})$。这也说明为何只对 $\lambda_i\in\mathbb{Z}$ 讨论：否则 $z^{\lambda_i}$ 不是单值幂级数，留数无定义。指代落点：§7.3.5，(7.3-15)--(7.3-18)，原书 p.217--219。[依据：原书 p.217--219][来源：本书 p.219。]}

% [Tutor Note: (7.3-16) 与 (7.3-17) 中 $p_{n+\lambda}$、$p_{i,n+\lambda_i}$ 的平移量 $\lambda$ 来自乘上的因子 $z^\lambda$：它在流的 Fourier 分解中把指标整体平移 $\lambda$。核对时别把 $\lambda$（权）与 $\lambda_i=\langle\lambda,\check\alpha_i\rangle$（该权的第 $i$ 个坐标）混为一谈。(7.3-18) 中求和下界是 $m\le\lambda_i$，与 (7.3-13) 的 $m\le0$ 在 $\lambda_i=0$ 时一致。]
